\section{AMC 12 12A Fall 2021}
        \begin{enumerate}
        	\item (\textit{AMC 12 12A Fall 2021}) What is the value of $\frac{(2112-2021)^2}{169}$?


 $\textbf{(A) } 7 \qquad\textbf{(B) } 21 \qquad\textbf{(C) } 49 \qquad\textbf{(D) } 64 \qquad\textbf{(E) } 91$

 \textbf{Answer: }C



\textbf{Solution 1 (Laws of Exponents)}

We have
 \[\frac{(2112-2021)^2}{169}=\frac{91^2}{169}=\frac{91^2}{13^2}=\left(\frac{91}{13}\right)^2=7^2=\boxed{\textbf{(C) } 49}.\]~MRENTHUSIASM

 



\textbf{Solution 2 (Difference of Squares)}

We have
 \[\frac{(2112-2021)^2}{169}=\frac{91^2}{169}=\frac{(10^2-3^2)^2}{13^2}=\frac{((10+3)(10-3))^2}{13^2}=\frac{(13\cdot7)^2}{13^2}=\frac{13^2 \cdot 7^2}{13^2}=7^2=\boxed{\textbf{(C) } 49}.\]




\textbf{Solution 3 (Estimate)}

We know that $2112-2021 = 91$. Approximate this as $100$ as it is pretty close to it. Also, approximate $169$ to $170$. We then have 
 \[\frac{(2112 - 2021)^2}{169} \approx \frac{100^2}{170} \approx \frac{1000}{17} \approx 58.\]Now check the answer choices. The two closest answers are $49$ and $64$. As the numerator is actually bigger than it should be, it should be the smaller answer, or $\boxed{\textbf{(C) } 49}$.

 


	\item (\textit{AMC 12 12A Fall 2021}) Menkara has a $4 \times 6$ index card. If she shortens the length of one side of this card by $1$ inch, the card would have area $18$ square inches. What would the area of the card be in square inches if instead she shortens the length of the other side by $1$ inch?


 $\textbf{(A) }16\qquad\textbf{(B) }17\qquad\textbf{(C) }18\qquad\textbf{(D) }19\qquad\textbf{(E) }20$

 \textbf{Answer: }E



\textbf{Solution}

We construct the following table:
 \[\begin{array}{c||c|c||c} & & & \\ [-2.5ex] \textbf{Scenario} & \textbf{Length} & \textbf{Width} & \textbf{Area} \\ [0.5ex] \hline & & & \\ [-2ex] \text{Initial} & 4 & 6 & 24 \\ \text{Menkara shortens one side.} & 3 & 6 & 18 \\ \text{Menkara shortens other side instead.} & 4 & 5 & 20 \end{array}\]Therefore, the answer is $\boxed{\textbf{(E) } 20}.$



	\item (\textit{AMC 12 12A Fall 2021}) Mr. Lopez has a choice of two routes to get to work. Route A is $6$ miles long, and his average speed along this route is $30$ miles per hour. Route B is $5$ miles long, and his average speed along this route is $40$ miles per hour, except for a $\frac{1}{2}$-mile stretch in a school zone where his average speed is $20$ miles per hour. By how many minutes is Route B quicker than Route A?


 $\textbf{(A)}\ 2 \frac{3}{4}  \qquad\textbf{(B)}\  3 \frac{3}{4} \qquad\textbf{(C)}\  4 \frac{1}{2} \qquad\textbf{(D)}\  5 \frac{1}{2} \qquad\textbf{(E)}\ 6 \frac{3}{4}$

 \textbf{Answer: }B



\textbf{Solution 1}

If Mr. Lopez chooses Route A, then he will spend $\frac{6}{30}=\frac{1}{5}$ hour, or $\frac{1}{5}\cdot60=12$ minutes.

 If Mr. Lopez chooses Route B, then he will spend $\frac{9/2}{40}+\frac{1/2}{20}=\frac{11}{80}$ hour, or $\frac{11}{80}\cdot60=8\frac14$ minutes.

 Therefore, Route B is quicker than Route A by $12-8\frac14=\boxed{\textbf{(B)}\  3 \frac{3}{4}}$ minutes.

 



\textbf{Solution 2}

We use the equation $d=st$ to solve this problem. Recall that $1$ mile per hour is equal to $\frac{1}{60}$ mile per minute.

 For Route A, the distance is $6$ miles and the speed to travel this distance is $\frac{1}{2}$ mile per minute. Thus, the time it takes on Route A is $12$ minutes.

 For Route B, we have to use the equation twice: once for the distance of $5- \frac{1}{2} = \frac{9}{2}$ miles with a speed of $\frac{2}{3}$ mile per minute and a distance of $\frac{1}{2}$ miles at a speed of $\frac{1}{3}$ mile per minute. Thus, the time it takes to go on Route B is $\frac{9}{2} \cdot \frac{3}{2} + \frac{1}{2} \cdot 3 = \frac{27}{4} + \frac{3}{2} = \frac{33}{4}$ minutes. Thus, Route B is $12 - \frac{33}{4} = \frac{15}{4} = 3\frac{3}{4}$ faster than Route A. Thus, the answer is $\boxed{\textbf{(B)}\  3 \frac{3}{4}}.$



	\item (\textit{AMC 12 12A Fall 2021}) The six-digit number $\underline{2}\,\underline{0}\,\underline{2}\,\underline{1}\,\underline{0}\,\underline{A}$ is prime for only one digit $A.$ What is $A?$
$(\textbf{A})\: 1\qquad(\textbf{B}) \: 3\qquad(\textbf{C}) \: 5 \qquad(\textbf{D}) \: 7\qquad(\textbf{E}) \: 9$

 \textbf{Answer: }E



\textbf{Solution 1}

First, modulo $2$ or $5$, $\underline{20210A} \equiv A$.Hence, $A \neq 0, 2, 4, 5, 6, 8$.

 Second modulo $3$, $\underline{20210A} \equiv 2 + 0 + 2 + 1 + 0 + A \equiv 5 + A$.Hence, $A \neq 1, 4, 7$.

 Third, modulo $11$, $\underline{20210A} \equiv A + 1 + 0 - 0 - 2 - 2 \equiv A - 3$.Hence, $A \neq 3$.

 Therefore, the answer is $\boxed{\textbf{(E)}\ 9}$.

 



\textbf{Solution 2 (Elimination)}

Any number ending in $5$ is divisible by $5$. So we can eliminate option $\textbf{(C)}$.

 If the sum of the digits of a number is divisible by $3$, the number is divisible by $3$. The sum of the digits of this number is $2 + 0 + 2 + 1 + 0 + A = 5 + A$. If $5 + A$ is divisible by $3$, the number is divisible by $3$. Thus we can eliminate options $\textbf{(A)}$ and $\textbf{(D)}$.

 So the correct option is either $\textbf{(B)}$ or $\textbf{(E)}$. Let's try dividing the number with some integers.

 $20210A/7 = 2887x$, where $x$ is $1A/7$. Since $13$ and $19$ are both indivisible by $7$, this does not help us narrow the choices down.

 $20210A/11 = 1837x$, where $x$ is $3A/11$. Since $33/11 = 3$, option $\textbf{(B)}$ would make $20210A$ divisible by $11$. Thus, by elimination, the correct choice must be option $\boxed{\textbf{(E)}\ 9}$.

 



\textbf{Solution 3}

$202100 \implies$ divisible by $2$.

 $202101 \implies$ divisible by $3$.

 $202102 \implies$ divisible by $2$.

 $202103 \implies$ divisible by $11$.

 $202104 \implies$ divisible by $2$.

 $202105 \implies$ divisible by $5$.

 $202106 \implies$ divisible by $2$.

 $202107 \implies$ divisible by $3$.

 $202108 \implies$ divisible by $2$.

 This leaves only $A=\boxed{\textbf{(E)}\ 9}$.

 \textbf{Sidenote} The divisibility test for $11$ is if the difference between the sum of the alternating digits is $0$ or $11$, then that number is divisible by $11$. For $202103$, we have $2$+$2$-$1$-$3$=$0$.

 For $7$, the divisibility test can be slow. Take the last digit, double it, and subtract it from the rest of the number. If it is divisible by $7$, then the original number is divisible by $7$. If the number is too large to determine whether or not it is a multiple of $7$, then we repeat the process until we can determine whether or not the number is divisible by $7$.

 


	\item (\textit{AMC 12 12A Fall 2021}) Elmer the emu takes $44$ equal strides to walk between consecutive telephone poles on a rural road. Oscar the ostrich can cover the same distance in $12$ equal leaps. The telephone poles are evenly spaced, and the $41$st pole along this road is exactly one mile ($5280$ feet) from the first pole. How much longer, in feet, is Oscar's leap than Elmer's stride?


 $\textbf{(A) }6\qquad\textbf{(B) }8\qquad\textbf{(C) }10\qquad\textbf{(D) }11\qquad\textbf{(E) }15$

 \textbf{Answer: }B



\textbf{Solution 1}

There are $41-1=40$ gaps between the $41$ telephone poles, so the distance of each gap is $5280\div40=132$ feet.

 Each of Oscar's leaps covers $132\div12=11$ feet, and each of Elmer's strides covers $132\div44=3$ feet.

 Therefore, Oscar's leap is $11-3=\boxed{\textbf{(B) }8}$ feet longer than Elmer's stride.

 



\textbf{Solution 2}

There are $41-1=40$ gaps between the $41$ telephone poles, so Elmer takes $44 \cdot 40 = 1760$ strides in total, and Oscar takes $12 \cdot 40 = 480$ leaps in total. Therefore, the answer is $(5280 \div 480) - (5280 \div 1760) = 11-3=\boxed{\textbf{(B) }8}$.

 


	\item (\textit{AMC 12 12A Fall 2021}) As shown in the figure below, point $E$ lies on the opposite half-plane determined by line $CD$ from point $A$ so that $\angle CDE = 110^\circ$. Point $F$ lies on $\overline{AD}$ so that $DE=DF$, and $ABCD$ is a square. What is the degree measure of $\angle AFE$?


 
                    \begin{figure}[H]
                        \centering
                        \includegraphics[width=0.5\textwidth]{Problem/AMC_12_12A_Fall_2021_6_Problem_Figure_1.png}
                    \end{figure}
                    
$\textbf{(A) }160\qquad\textbf{(B) }164\qquad\textbf{(C) }166\qquad\textbf{(D) }170\qquad\textbf{(E) }174$

 \textbf{Answer: }D



\textbf{Solution 1}

By angle subtraction, we have $\angle ADE = 360^\circ - \angle ADC - \angle CDE = 160^\circ.$ Note that $\triangle DEF$ is isosceles, so $\angle DFE = \frac{180^\circ - \angle ADE}{2}=10^\circ.$ Finally, we get $\angle AFE = 180^\circ - \angle DFE = \boxed{\textbf{(D) }170}$ degrees.

 



\textbf{Solution 2 (Extension)}

We can extend $\overline{AD}$ to $G$, making $\angle CDG$ a right angle. It follows that $\angle GDE$ is $110^\circ - 90^\circ = 20^\circ$, as shown below.
                    \begin{figure}[H]
                        \centering
                        \includegraphics[width=0.5\textwidth]{Solution/AMC_12_12A_Fall_2021_6_Solution_Figure_1.png}
                    \end{figure}
                    Since $\angle DFE = \angle DEF$, we see that $\angle DFE = \angle DEF = \frac{20}{2} = 10^\circ$. Thus, $\angle AFE = 180^\circ - 10^\circ = \boxed{\textbf{(D)} ~170}$ degrees.

 


	\item (\textit{AMC 12 12A Fall 2021}) A school has $100$ students and $5$ teachers. In the first period, each student is taking one class, and each teacher is teaching one class. The enrollments in the classes are $50, 20, 20, 5,$ and $5$. Let $t$ be the average value obtained if a teacher is picked at random and the number of students in their class is noted. Let $s$ be the average value obtained if a student was picked at random and the number of students in their class, including the student, is noted. What is $t-s$?


 $\textbf{(A)}\ {-}18.5  \qquad\textbf{(B)}\ {-}13.5 \qquad\textbf{(C)}\ 0 \qquad\textbf{(D)}\ 13.5 \qquad\textbf{(E)}\ 18.5$

 \textbf{Answer: }B



\textbf{Solution}

The formula for expected values is 
 \[\text{Expected Value}=\sum(\text{Outcome}\cdot\text{Probability}).\]We have
 \begin{align*} t &= 50\cdot\frac15 + 20\cdot\frac15 + 20\cdot\frac15 + 5\cdot\frac15 + 5\cdot\frac15 \\ &= (50+20+20+5+5)\cdot\frac15 \\ &= 100\cdot\frac15 \\ &= 20, \\ s &= 50\cdot\frac{50}{100} + 20\cdot\frac{20}{100} + 20\cdot\frac{20}{100} + 5\cdot\frac{5}{100} + 5\cdot\frac{5}{100} \\ &= 25 + 4 + 4 + 0.25 + 0.25 \\ &= 33.5. \end{align*}
Therefore, the answer is $t-s=\boxed{\textbf{(B)}\ {-}13.5}.$



	\item (\textit{AMC 12 12A Fall 2021}) Let $M$ be the least common multiple of all the integers $10$ through $30,$ inclusive. Let $N$ be the least common multiple of $M,32,33,34,35,36,37,38,39,$ and $40.$ What is the value of $\frac{N}{M}?$
$\textbf{(A)}\ 1 \qquad\textbf{(B)}\ 2 \qquad\textbf{(C)}\ 37 \qquad\textbf{(D)}\ 74 \qquad\textbf{(E)}\ 2886$

 \textbf{Answer: }D



\textbf{Solution}

By the definition of least common mutiple, we take the greatest powers of the prime numbers of the prime factorization of all the numbers, that we are taking the $\text{lcm}$ of. In this case, 
 \[M = 2^4 \cdot 3^3 \cdot 5^2 \cdot 7 \cdot 11 \cdot 13 \cdot 17 \cdot 19 \cdot 23 \cdot 29.\] Now, using the same logic, we find that 
 \[N = M \cdot 2 \cdot 37,\] because we have an extra power of $2$ and an extra power of $37.$ Thus, $\frac{N}{M} = 2\cdot 37 = \boxed{\textbf{(D)}\ 74}.$




\textbf{Solution 2}

We know that between 32 and 40, there is one prime number, namely $37$. This was not included in the previous $M$, so we know that the value of $\frac{N}{M}$ will include a new 37. Then, let’s consider each individual number. 
 \[32 = 2^5\] has a new power of 2 now included in $M$, so we add that to our list. 
 \[33=3 \cdot 11\], which was included in the previous $M$. Similarly, we can go through the numbers 32 to 40 and find that the only new powers of numbers added are 
 \[2 \cdot 37 = 74\], which gives $\frac{N}{M} = 2\cdot 37 = \boxed{\textbf{(D)}\ 74}.$



	\item (\textit{AMC 12 12A Fall 2021}) A right rectangular prism whose surface area and volume are numerically equal has edge lengths $\log_{2}x, \log_{3}x,$ and $\log_{4}x.$ What is $x?$
$\textbf{(A)}\ 2\sqrt{6} \qquad\textbf{(B)}\ 6\sqrt{6} \qquad\textbf{(C)}\ 24 \qquad\textbf{(D)}\ 48 \qquad\textbf{(E)}\ 576$

 \textbf{Answer: }E



\textbf{Solution 1}

The surface area of this right rectangular prism is $2(\log_{2}x\log_{3}x+\log_{2}x\log_{4}x+\log_{3}x\log_{4}x).$

 The volume of this right rectangular prism is $\log_{2}x\log_{3}x\log_{4}x.$

 Equating the numerical values of the surface area and the volume, we have 
 \[2(\log_{2}x\log_{3}x+\log_{2}x\log_{4}x+\log_{3}x\log_{4}x)=\log_{2}x\log_{3}x\log_{4}x.\]Dividing both sides by $\log_{2}x\log_{3}x\log_{4}x,$ we get 
 \[2\left(\frac{1}{\log_{4}x}+\frac{1}{\log_{3}x}+\frac{1}{\log_{2}x}\right)=1. \hspace{15mm} (\bigstar)\]Recall that $\log_{b}a=\frac{1}{\log_{a}b}$ and $\log_{b}\left(a^n\right)=n\log_{b}a,$ so we rewrite $(\bigstar)$ as 
 \begin{align*} 2(\log_{x}4+\log_{x}3+\log_{x}2)&=1 \\ 2\log_{x}24&=1 \\ \log_{x}576&=1 \\ x&=\boxed{\textbf{(E)}\ 576}. \end{align*}
~MRENTHUSIASM

 



\textbf{Solution 2}

We know that the formulae for the surface area and volume of a right rectangular prism are $S=2lw+2wh+2lh$ and $V=lwh$, respectively.It is given that these values are numerically equal, so we can set $2lw+2wh+2lh=lwh$.Dividing both sides by $lwh$ gives$\frac{2}{h} + \frac{2}{l} + \frac{2}{w} = 1$.We then substitute values for $l$, $w$, and $h$, which yields $\frac{2}{\log_2 x} + \frac{2}{\log_3 x} + \frac{2}{\log_4 x}$.Then, we can recall the logarithm base-change formula $\log_b a = \frac{\log_x a}{\log_x b}$. We can use this formula to change all the sidelengths to log base $2$ as follows:
 \[\log_3 x = \frac{\log_2 x}{\log_2 3}\]
 \[\log_4 x = \frac{\log_2 x}{\log_2 4}\]Then substitute back into our equation to get
 \[1 = \frac{2}{\log_2 x} + \frac{2}{\dfrac{\log_2 x}{\log_2 3}} + \frac{2}{\dfrac{\log_2 x}{\log_2 4}}\].Simplifying this equation, as well as the fact that $\log_2 4 = 2$, we get:
 \[1 = \frac{2}{\log_2 x} + \frac{2\log_2 3}{\log_2 x} + \frac{4}{\log_2 x}\]Further simplification yields:\begin{align*}1 &= \frac{2+2\log_2 3+4}{\log_2 x} \\1 &= \frac{6+\log_2 9}{\log_2 x} \\\log_2 x &= 6+\log_2 9 \\\log_2 x - \log_2 9 &= 6 \\\log_2 \frac{x}{9} &= 6 \\\frac{x}{9} &= 2^6 \\x &= 9 \cdot 64 = \boxed{\text{(E) } 576}\end{align*}


 


	\item (\textit{AMC 12 12A Fall 2021}) The base-nine representation of the number $N$ is $27{,}006{,}000{,}052_{\text{9}}.$ What is the remainder when $N$ is divided by $5?$
$\textbf{(A) } 0\qquad\textbf{(B) } 1\qquad\textbf{(C) } 2\qquad\textbf{(D) } 3\qquad\textbf{(E) }4$

 \textbf{Answer: }D



\textbf{Solution 1 (Modular Arithmetic)}

Recall that $9\equiv-1\pmod{5}.$ We expand $N$ by the definition of bases:
 \begin{align*} N&=27{,}006{,}000{,}052_9 \\ &= 2\cdot9^{10} + 7\cdot9^9 + 6\cdot9^6 + 5\cdot9 + 2 \\ &\equiv 2\cdot(-1)^{10} + 7\cdot(-1)^9 + 6\cdot(-1)^6 + 2 &&\pmod{5} \\ &\equiv 2-7+6+2 &&\pmod{5} \\ &\equiv \boxed{\textbf{(D) } 3} &&\pmod{5}. \end{align*}
~Aidensharp ~Kante314 ~MRENTHUSIASM ~anabel.disher

 



\textbf{Solution 2 (Powers of 9)}

We need to first convert $N$ into a regular base-$10$ number: 
 \[N = 27{,}006{,}000{,}052_9 = 2\cdot9^{10} + 7\cdot9^9 + 6\cdot9^6 + 5\cdot9 + 2.\]

 Now, consider how the last digit of $9$ changes with changes of the power of $9:$
 \begin{align*} 9^0&=1, \\ 9^1&=9, \\ 9^2&=\ldots 1, \\ 9^3&=\ldots 9, \\ 9^4&=\ldots 1, \\ & \ \vdots \end{align*}
Note that if $x$ is odd, then $9^x \equiv 4\pmod{5}.$ On the other hand, if $x$ is even, then $9^x \equiv 1\pmod{5}.$

 Therefore, we have
 \begin{align*} N&\equiv 2\cdot(1) + 7\cdot(4) + 6\cdot(1) + 5\cdot(4) + 2\cdot(1) &&\pmod{5} \\ &\equiv 2+28+6+20+2 &&\pmod{5} \\ &\equiv 58 &&\pmod{5} \\ &\equiv \boxed{\textbf{(D) } 3} &&\pmod{5}. \\ \end{align*}
Note that for the odd case, $9^x \equiv -1\pmod{5}$ may simplify the process further, as given by Solution 1.

 \textbf{Remark} By the long division algorithm, you can work from left to right accumulating the answer, and don't need to count the digits to get the power. You only need to take the current leading digit modulo $5$, negate it, and add it to the next digit, and repeat the process until you reach the units digit.

 


	\item (\textit{AMC 12 12A Fall 2021}) Consider two concentric circles of radius $17$ and $19.$ The larger circle has a chord, half of which lies inside the smaller circle. What is the length of the chord in the larger circle?


 $\textbf{(A)}\ 12\sqrt{2} \qquad\textbf{(B)}\ 10\sqrt{3} \qquad\textbf{(C)}\ \sqrt{17 \cdot 19} \qquad\textbf{(D)}\ 18 \qquad\textbf{(E)}\ 8\sqrt{6}$

 \textbf{Answer: }E



\textbf{Solution 1 (Pythagorean Theorem)}

Label the center of both circles $O$. Label the chord in the larger circle as $\overline{ABCD}$, where $A$ and $D$ are on the larger circle and $B$ and $C$ are on the smaller circle. Construct the radius perpendicular to the chord and label their intersection as $M$. Because a radius that is perpendicular to a chord bisects the chord, $M$ is the midpoint of the chord. 

 Construct segments $\overline{AO}$ and $\overline{BO}$. These are radii with lengths 17 and 19 respectively. 

 Then, use the \href{https://artofproblemsolving.com/wiki/index.php?title=Pythagorean\_Theorem}{Pythagorean Theorem}. In $\triangle OMA$, we have
 \begin{align*} OM^2 & = OA^2 - AM^2 \\ & = OA^2 - \left( \frac{AD}{2} \right)^2 \\ & = 19^2 - \frac{AD^2}{4} . \end{align*}


 In $\triangle OMB$, we have
 \begin{align*} OM^2 & = OB^2 - BM^2 \\ & = OB^2 - \left( \frac{BC}{2} \right)^2 \\ & = OB^2 - \left( \frac{AD}{4} \right)^2 \\ & = 17^2 - \frac{AD^2}{16} . \end{align*}


 Equating these two expressions, we get 
 \[19^2 - \frac{AD^2}{4} = 17^2 - \frac{AD^2}{16}\]and $AD=\boxed{\textbf{(E) }8 \sqrt{6}}$.

 



\textbf{Solution 2 (Power of a Point)}

Draw the diameter perpendicular to the chord. Notice that by symmetry this diameter bisects the chord.

 Call the intersection between that diameter and the chord $A$. In the smaller circle, let the shorter piece of the diameter cut by the chord be $x$, making the longer piece $34-x.$ In that same circle, let the $y$ be the length of the portion of the chord that is cut by the diameter. 

 The radius of the larger circle is $2$ more than the radius of the small circle. So, in the larger circle, the shorter piece of the diameter cut by the chord is of length $x+2$ and the longer piece is $36-x.$ As given in the problem, the bisected length of the chord in the larger circle is twice as much, so it is of length $2y$. By \href{https://artofproblemsolving.com/wiki/index.php?title=Power\_of\_a\_Point\_Theorem}{Power of a Point}, we can construct a system of equations
 \[x(34-x) = y^2\]
 \[(x+2)(36-x) = (2y)^2.\]Expanding both equations, we get
 \[34x-x^2 = y^2\]
 \[36x-x^2+72-2x = 4y^2,\]in which the $34x$ and $-x^2$ terms magically cancel when we subtract the first equation from the second equation. Thus, now we have 
 \[72 = 3y^2 \implies y = 2\sqrt{6} \implies 4y = \boxed{\textbf{(E)} \: 8\sqrt{6}}.\]

 -fidgetboss\_4000

 



\textbf{Solution 3 (Coordinate Geometry)}

Represent the circles as $x^{2}+y^{2}=17^{2}$, and $x^{2}+y^{2}=19^{2}$. Solving for $x$ in these equations we obtain $x=\sqrt{17^{2}-y^{2}}$ and $x=\sqrt{19^{2}-y^{2}}$. Because half of the chord is in the smaller circle, the larger circle should have an $x$ value that is twice as big as the smaller circle's $x$ value. We set up the equation:
 \[2\sqrt{289-y^{2}}=\sqrt{361-y^{2}}\] 
 \[4(289-y^{2})=361-y^{2}\]
 \[y^{2}=265\]Substituting $y^{2}$ into $x^{2}+y^{2}=19^{2}$, we obtain $x=\sqrt{96}=4\sqrt{6}$. However, this is only half of the chord length, so we must double it to obtain $\boxed{\textbf{(E)} \: 8\sqrt{6}}.$

 -helpmebro

 


 


	\item (\textit{AMC 12 12A Fall 2021}) What is the number of terms with rational coefficients among the $1001$ terms in the expansion of $\left(x\sqrt[3]{2}+y\sqrt{3}\right)^{1000}?$
$\textbf{(A)}\ 0 \qquad\textbf{(B)}\ 166 \qquad\textbf{(C)}\ 167 \qquad\textbf{(D)}\ 500 \qquad\textbf{(E)}\ 501$

 \textbf{Answer: }C



\textbf{Solution 1}

By the Binomial Theorem, each term in the expansion is of the form 
 \[\binom{1000}{k}\left(x\sqrt[3]{2}\right)^k\left(y\sqrt{3}\right)^{1000-k}=\binom{1000}{k}2^{\frac k3}3^{\frac{1000-k}{2}}x^k y^{1000-k},\] where $k\in\{0,1,2,\ldots,1000\}.$

 This problem is equivalent to counting the values of $k$ such that both $\frac k3$ and $\frac{1000-k}{2}$ are integers. Note that $k$ must be a multiple of $3$ and a multiple of $2,$ so $k$ must be a multiple of $6.$ There are $\boxed{\textbf{(C)}\ 167}$ such values of $k:$ 
 \[6(0), 6(1), 6(2), \ldots, 6(166).\]




\textbf{Solution 2 (Rough Conceptual Approach)}

Similar to Solution 1, however, instead of applying the Binomial Theorem, we can notice that it "takes" two $\;\sqrt{3}\;$ terms multiplied by each other to get a rational coefficient for $y$ and it "takes" three $\;\sqrt[3]{2}\;$ terms multiplied by each other to get a rational coefficient for $x$. Since the only way to get a total rational coefficient (that results when you multiply the final coefficient of $y$ by that of for $x$) is for both $x$'s and $y$'s coefficients to be rational, we can notice that the solution simply comes down to the number of multiples of $6$ between $0$ and $1000$, which turns out to be $\boxed{\textbf{(C)}\ 167}$.

 


	\item (\textit{AMC 12 12A Fall 2021}) The angle bisector of the acute angle formed at the origin by the graphs of the lines $y = x$ and $y=3x$ has equation $y=kx.$ What is $k?$
$\textbf{(A)} \ \frac{1+\sqrt{5}}{2} \qquad \textbf{(B)} \ \frac{1+\sqrt{7}}{2} \qquad \textbf{(C)} \ \frac{2+\sqrt{3}}{2} \qquad \textbf{(D)} \ 2\qquad \textbf{(E)} \ \frac{2+\sqrt{5}}{2}$

 \textbf{Answer: }A



\textbf{Solution 1 (Angle Bisector Theorem)}

This solution refers to the \textbf{Diagram} section.

 Let $O=(0,0), A=(3,3), B=(1,3),$ and $C=\left(\frac3k,3\right).$ As shown below, note that $\overline{OA}, \overline{OB},$ and $\overline{OC}$ are on the lines $y=x, y=3x,$ and $y=kx,$ respectively. By the Distance Formula, we have $OA=3\sqrt2, OB=\sqrt{10}, AC=3-\frac3k,$ and $BC=\frac3k-1.$
                    \begin{figure}[H]
                        \centering
                        \includegraphics[width=0.5\textwidth]{Solution/AMC_12_12A_Fall_2021_13_Solution_Figure_1.png}
                    \end{figure}
                    By the Angle Bisector Theorem, we get $\frac{OA}{OB}=\frac{AC}{BC},$ or 
 \begin{align*} \frac{3\sqrt2}{\sqrt{10}}&=\frac{3-\frac3k}{\frac3k-1} \\ \frac{3\sqrt2}{\sqrt{10}}&=\frac{3k-3}{3-k} \\ \frac{\sqrt2}{\sqrt{10}}&=\frac{k-1}{3-k} \\ \frac15&=\frac{(k-1)^2}{(3-k)^2} \\ 5(k-1)^2&=(3-k)^2 \\ 4k^2-4k-4&=0 \\ k^2-k-1&=0 \\ k&=\frac{1\pm\sqrt5}{2}. \end{align*}
Since $k>0,$ the answer is $k=\boxed{\textbf{(A)} \ \frac{1+\sqrt{5}}{2}}.$
<u>\textbf{Remark}</u>

 The value of $k$ is known as the \textbf{Golden Ratio}: $\phi=\frac{1+\sqrt{5}}{2}\approx 1.61803398875.$




\textbf{Solution 2 (Analytic and Plane Geometry)}


                    \begin{figure}[H]
                        \centering
                        \includegraphics[width=0.5\textwidth]{Solution/AMC_12_12A_Fall_2021_13_Solution_Figure_2.png}
                    \end{figure}
                    Consider the graphs of $f(x)=x$ and $g(x)=3x$. Since it will be easier to consider at unity, let $x=1$, then we have $f(1)=1$ and $g(1)=3$.

 Now, let $O$ be $(0,0)$, $A$ be $(1,1)$, and $B$ be $(1,3)$. Cutting through side $AB$ of triangle $OAB$ is the angle bisector $OC$ where $C$ is on side $AB$.

 Hence, by the Angle Bisector Theorem, we get $\frac{OB}{OA}=\frac{BC}{AC}$. 

 By the Pythagorean Theorem, $OA=\sqrt{2}$ and $OB=\sqrt{10}$. Therefore, $\frac{BC}{AC}=\sqrt{5} \implies BC=\sqrt{5}AC$.

 Since $AB=AC+BC=2$, it is easy derive $AC+\sqrt{5}AC=2 \implies AC=\frac{2}{1+\sqrt{5}}=\frac{-1+\sqrt{5}}{2}$.

 The vertical distance between the $x$-axis and $C$ is $\frac{-1+\sqrt{5}}{2}+1=\frac{1+\sqrt{5}}{2}$. Because the $x$-coordinate of point $C$ is $1$, the slope we need to find is just the $y$-coordinate $\boxed{\textbf{(A)} \ \frac{1+\sqrt{5}}{2}}.$




\textbf{Solution 3 (Analytic and Plane Geometry)}

Let's begin by drawing a triangle that starts at the origin. Assume that the base of the triangle goes to the point $x = 1$. The line $x = y$ is the hypotenuse of a right triangle with side length $1$. The hypotenuse' length is $\sqrt 2$. Then, let's draw the line $x = 3y$. We extend it to when $x = 1$. The length of the hypotenuse of the larger triangle is $\sqrt {10}$ with legs $1, 3$. We then draw the angle bisector. We should label the triangle, so here we go. $AC$ is $1$. $BC$ is $3$. $AB$ is $\sqrt {10}$. When the line with angle $45 ^\circ$ intersects the line $x = 1$, call the point $D$. When the angle bisector intersects the line $x = 1$, call the point $E$. By Angle Bisector Theorem, $\frac {DE}{DB} = \frac {\sqrt {2}}{\sqrt{10}}$. Since $BC$ is $3$ and $DC$ is $1$, we have that $BD$ is $2$. Solving for $DE$, we get that $DE$ is $\frac {\sqrt 5 - 1}{2}$. 

 Since $DE$ is $\frac {\sqrt 5 - 1}{2}$, we have that $CE$ is just one more than that. Therefore, $CE$ is $\frac {1+\sqrt 5}{2}$. Since $AC$ is $1$, we get that $k$ is $\boxed{\textbf{(A)} \ \frac{1+\sqrt{5}}{2}}$.

 <u>\textbf{Remark}</u>

 The answer turns out to be the golden ratio or phi ($\phi$). Phi has many properties and is related to the \href{https://artofproblemsolving.com/wiki/index.php?title=Fibonacci\_sequence}{Fibonacci sequence}. See \href{https://artofproblemsolving.com/wiki/index.php?title=Phi}{Phi}.

 



\textbf{Solution 4 (Distance Between a Point and a Line)}

Note that the distance between the point $(m,n)$ to line $Ax + By + C = 0,$ is $\frac{|Am + Bn +C|}{\sqrt{A^2 +B^2}}.$ Because line $y=kx$ is a perpendicular bisector, a point on the line $y=kx$ must be equidistant from the two lines($y=x$ and $y=3x$), call this point $P(z,w).$ Because, the line $y=kx$ passes through the origin, our requested value of $k,$ which is the slope of the angle bisector line, can be found when evaluating the value of $\frac{w}{z}.$ By the Distance from Point to Line formula we get the equation, 
 \[\frac{|3z-w|}{\sqrt{10}} = \frac{|z-w|}{\sqrt{2}}.\] Note that $|3z-w|\ge 0,$ because $y=3x$ is higher than $P$ and $|z-w|\le 0,$ because $y=x$ is lower to $P.$ Thus, we solve the equation, 
 \[(3z-w)\sqrt{2} = (w-z)\sqrt{10} \Rightarrow  3z-w = \sqrt{5} \cdot(w-z)\Rightarrow (\sqrt{5} +1)w = (3+\sqrt{5})z.\] Thus, the value of $\frac{w}{z} = \frac{3+\sqrt{5}}{1+\sqrt{5}} = \frac{1+\sqrt{5}}{2}.$ Thus, the answer is $\boxed{\textbf{(A)} \ \frac{1+\sqrt{5}}{2}}.$

 (Fun Fact: The value $\frac{1+\sqrt{5}}{2}$ is the golden ratio $\phi.$)

 



\textbf{Solution 5 (Trigonometry)}

This problem can be trivialized using basic trig identities. Let the angle made by $y=x$ and the $x$-axis be $\theta_{1}$ and the angle made by $y=3x$ and the $x$-axis be $\theta_{3}$. Note that $\tan(\theta_{1})=1$ and $\tan(\theta_{3})=3$, and this is why we named them as such. Let the angle made by $y=kx$ be denoted as $\theta_{k}$. Since $y=kx$ bisects the two lines, notice that
 \[\theta_k-\theta_1=\theta_3-\theta_k.\]

 Now, we can take the tangent and apply the tangent subtraction formula:
 \begin{align*} \tan(\theta_k-\theta_1)&=\tan(\theta_3-\theta_k)\\ \frac{\tan(\theta_k)-\tan(\theta_1)}{1+\tan(\theta_k)\tan(\theta_1)}&=\frac{\tan(\theta_3)-\tan(\theta_k)}{1+\tan(\theta_3)\tan(\theta_k)}\\ \frac{k-1}{1+k}&=\frac{3-k}{1+3k}\\ (k-1)(1+3k)&=(1+k)(3-k)\\ 3k^2-2k-1&=-k^2+2k+3\\ 4k^2-4k-4&=0\\ k^2-k-1&=0\\ \implies k&=\frac{1\pm \sqrt{5}}{2} \end{align*}
Since $y=kx$ is increasing, $k>0$; thus, $k=\boxed{\textbf{(A)} \ \frac{1+\sqrt{5}}{2}}.$




\textbf{Solution 6 (Trigonometry)}

Denote by $\alpha_1$, $\alpha_2$, $\alpha_3$ the acute angles formed between the $x$-axis and lines $y = x$, $y = 3 x$, $y = k x$, respectively.Hence, $\tan \alpha_1 = 1$, $\tan \alpha_2 = 3$, $\tan \alpha_3 = k$.

 Denote by $\theta$ the acute angle formed by lines $y = x$ and $y = 3 x$.

 Hence,
 \begin{align*} \tan \theta & = \tan \left( \alpha_2 - \alpha_1 \right) \\ & = \frac{\tan \alpha_2 - \tan \alpha_1}{1 + \tan \alpha_1 \tan \alpha_2} \\ &= \frac{3 - 1}{1 + 1 \cdot 3} \\ & = \frac{1}{2} . \end{align*}


 Following from the double-angle identity, we have
 \[ \tan \theta = \frac{2 \tan \frac{\theta}{2}}{1 - \tan^2 \frac{\theta}{2}} . \]

 Hence, $\tan \frac{\theta}{2} = - 2 \pm \sqrt{5}$.

 Because $\theta$ is acute, $\frac{\theta}{2}$ is acute.Hence, $\tan \frac{\theta}{2} > 0$.Hence, $\tan \frac{\theta}{2} = - 2 + \sqrt{5}$.

 Because line $y = kx$ is the angle bisector of $\theta$, the angle between lines $y = x$ and $y = k x$ is $\frac{\theta}{2}$.

 Hence,
 \begin{align*} \tan \alpha_3 & = \tan \left( \alpha_1 + \frac{\theta}{2} \right) \\ & = \frac{\tan \alpha_1 + \tan \frac{\theta}{2}}{1 - \tan \alpha_1 \tan \frac{\theta}{2}} \\ &= \frac{1 + \left( - 2 + \sqrt{5} \right)}{1 - 1 \cdot \left( - 2 + \sqrt{5} \right) } \\ & = \frac{\sqrt{5} - 1}{3 - \sqrt{5}} \\ & = \frac{1 + \sqrt{5}}{2} . \end{align*}


 Therefore, the answer is $\boxed{\textbf{(A) }\frac{1 + \sqrt{5}}{2}}$.

 



\textbf{Solution 7 (Vectors)}


                    \begin{figure}[H]
                        \centering
                        \includegraphics[width=0.5\textwidth]{Solution/AMC_12_12A_Fall_2021_13_Solution_Figure_3.png}
                    \end{figure}
                    When drawing the lines $y=x$ and $y=3x$, it is natural to choose points $A(1,1)$ and $B(1,3)$ together with origin $O$. See the figure attached. We utilize the fact that if $\mathbf{u}$ and $\mathbf{v}$ are vectors of same length, then $\mathbf{u} + \mathbf{v}$ bisects the angle between $\mathbf{u}$ and $\mathbf{v}$.

 In particular, we scale the vector $\overrightarrow{OA} = (1,1)$ by the factor of $\frac{OB}{OA} = \frac{\sqrt{10}}{\sqrt{2}} = \sqrt{5}$ to get $\overrightarrow{OA'} \coloneqq \sqrt{5}\,\overrightarrow{OA} = \left(\sqrt{5}, \sqrt{5}\right)$.So by adding vectors $\overrightarrow{OA'}$ and $\overrightarrow{OB} = (1,3)$ we get
 \[{\color[rgb]{0.666667,0,0}\overrightarrow{OC}}      \coloneqq {\color[rgb]{0,0.4,0.65}\overrightarrow{OA'}} + {\color[rgb]{0,0.4,0.65}\overrightarrow{OB}}      = \left( 1 + \sqrt{5}, 3 + \sqrt{5} \right)\]which bisects the acute angle formed by lines $OA: y = x$ and $OB: y = 3x$. (In other words, quadrilateral $OBCA'$ is a rhombus.)Finally, observe that $C\!\left(1+\sqrt{5}, 3+\sqrt{5}\right)$ lies on the line $y = kx$ whose slope is
 \[k = \frac{3+\sqrt{5}}{1+\sqrt{5}} = \frac{1+\sqrt{5}}{2}.\]Thus, the answer is $\boxed{\textbf{(A)}\;\frac{1+\sqrt{5}}{2}}$. $\blacksquare$




\textbf{Solution 8 (Dot Product)}

We notice that the line $y = x$ can be represented as the vector $\vec{A} = \begin{bmatrix}      1 \\ 1     \end{bmatrix}$and $y = 3x$ as $\vec{B} = \begin{bmatrix}      1 \\ 3     \end{bmatrix}$as the "slope" of both vectors represent the coefficient of $x$.

 Then, we can represent $y = kx$ as $\vec{C} = \begin{bmatrix}      1 \\ k     \end{bmatrix}$ and notice that since $C$ is in essence an angle bisector, 
 \[\theta = \angle \vec{A}\vec{C} = \angle \vec{B}\vec{C}\]

 \[A \cdot C = |A|\cdot|C|\cos(\theta_1)\] where $\theta_1 = \angle \vec{A}\vec{C} = \theta$
 \[B \cdot C = |B|\cdot|C|\cos(\theta_2)\] where $\theta_2 = \angle \vec{B}\vec{C} = \theta$

 Since both $\theta_i$'s are equivalent, we may simply represent them with $\theta$.

 Simplifying both equations by performing the necessary operations, we get

 
 \[1+k = \sqrt{2} \cdot \sqrt{k^2 + 1} \cdot \cos(\theta) \implies \frac{1+k}{\sqrt{2}} = \sqrt{k^2 + 1} \cos(\theta)\]

 \[1+3k = \sqrt{10} \cdot \sqrt{k^2 + 1} \cdot \cos(\theta)\]

 Substituting the first into the second, we get
 \[1 + 3k = \sqrt{10} \cdot \frac{1+k}{\sqrt{2}} \implies k=\boxed{\textbf{(A)} \ \frac{1+\sqrt{5}}{2}}\]




\textbf{Solution 9 (Trig by Averaging the Angles)}

Let $\theta_1$ and $\theta_2$ be the angles that the lines $y = x$ and $y = 3x$ make with the positive $x$-axis, so $\tan \theta_1 = 1$ and $\tan \theta_2 = 3$. 

 The slope $k$ of the angle bisector is given by $k = \tan \theta$, where $\theta = \frac{\theta_1 + \theta_2}{2}$. Using the tangent addition formula, we find $\tan(2\theta)$:

 
 \[tan(2\theta) = \tan(\theta_1 + \theta_2) = \frac{\tan \theta_1 + \tan \theta_2}{1 - \tan \theta_1 \tan \theta_2} = \frac{1 + 3}{1 - (1 \cdot 3)} = -2\]Applying the double-angle identity $\tan(2\theta) = \frac{2\tan\theta}{1 - \tan^2\theta}$ and substituting $k = \tan\theta$:
 \[\frac{2k}{1 - k^2} = -2 \implies -k = 1 - k^2 \implies k^2 - k - 1 = 0\]Solving this quadratic equation yields $k = \frac{1 \pm \sqrt{5}}{2}$. Since both lines lie in the first quadrant, their acute angle bisector must have a positive slope. Thus, we select the positive root:
 \[k = \frac{1+\sqrt{5}}{2}\]



	\item (\textit{AMC 12 12A Fall 2021}) In the figure, equilateral hexagon $ABCDEF$ has three nonadjacent acute interior angles that each measure $30^\circ$. The enclosed area of the hexagon is $6\sqrt{3}$. What is the perimeter of the hexagon?

                    \begin{figure}[H]
                        \centering
                        \includegraphics[width=0.5\textwidth]{Problem/AMC_12_12A_Fall_2021_14_Problem_Figure_1.png}
                    \end{figure}
                    
$\textbf{(A)} \: 4 \qquad \textbf{(B)} \: 4\sqrt3 \qquad \textbf{(C)} \: 12 \qquad \textbf{(D)} \: 18 \qquad \textbf{(E)} \: 12\sqrt3$

 \textbf{Answer: }E



\textbf{Solution 1}

Divide the equilateral hexagon $ABCDEF$ into isosceles triangles $ABF$, $CBD$, and $EDF$ and triangle $BDF$. The three isosceles triangles are congruent by SAS congruence. By \href{https://artofproblemsolving.com/wiki/index.php?title=CPCTC}{CPCTC}, $BF=BD=DF$, so triangle $BDF$ is equilateral. 

 Let the side length of the hexagon be $s$. The area of each isosceles triangle is 
 \[\frac{1}{2} a b \sin\angle C = \frac{1}{2} \cdot s \cdot s \cdot \sin{30^{\circ}} = \frac{1}{4}s^2.\]

 By the \href{https://artofproblemsolving.com/wiki/index.php?title=Law\_of\_Cosines}{Law of Cosines} on triangle $ABF$, 
 \[BF^2=s^2+s^2-2s^2\cos{30^{\circ}}=2s^2-\sqrt{3}s^2.\] 

 Hence, the 

 


	\item (\textit{AMC 12 12A Fall 2021}) Recall that the conjugate of the complex number $w = a + bi$, where $a$ and $b$ are real numbers and $i = \sqrt{-1}$, is the complex number $\overline{w} = a - bi$. For any complex number $z$, let $f(z) = 4i\hspace{1pt}\overline{z}$. The polynomial 
 \[P(z) = z^4 + 4z^3 + 3z^2 + 2z + 1\] has four complex roots: $z_1$, $z_2$, $z_3$, and $z_4$. Let 
 \[Q(z) = z^4 + Az^3 + Bz^2 + Cz + D\] be the polynomial whose roots are $f(z_1)$, $f(z_2)$, $f(z_3)$, and $f(z_4)$, where the coefficients $A,$ $B,$ $C,$ and $D$ are complex numbers. What is $B + D?$
$(\textbf{A})\: {-}304\qquad(\textbf{B}) \: {-}208\qquad(\textbf{C}) \: 12i\qquad(\textbf{D}) \: 208\qquad(\textbf{E}) \: 304$

 \textbf{Answer: }D



\textbf{Solution 1}

By Vieta's formulas, $z_1z_2z_3z_4=1$, and $D= (4i)^4\overline{z}_1\,\overline{z}_2\,\overline{z}_3\,\overline{z}_4.$

 Since $\overline{a}\cdot\overline{b}=\overline{ab},$
 \[D=(4i)^4\overline{z_1z_2z_3z_4} = 256(\overline{1}) = 256\]

 By Vieta's formulas, $z_1z_2+z_1z_3+\dots+z_3z_4=3$, and $B=(4i)^2\left(\overline{z}_1\,\overline{z}_2+\overline{z}_1\,\overline{z}_3+\dots+\overline{z}_3\,\overline{z}_4\right).$

 Since $\overline{a}\cdot\overline{b}=\overline{ab},$ 
 \[B=(4i)^2\left(\overline{z_1z_2}+\overline{z_1z_3}+\overline{z_1z_4}+\overline{z_2z_3}+\overline{z_2z_4}+\overline{z_3z_4}\right).\]Since $\overline{a}+\overline{b}=\overline{a+b},$
 \[B=(4i)^2\left(\overline{z_1z_2+z_1z_3+\dots+z_3z_4}\right)=-16(\overline{3})=-48\]

 
Our answer is $B+D=256-48=\boxed{(\textbf{D}) \: 208}.$




\textbf{Solution 2}

Since the coefficients of $P$ are real, the roots of $P$ can also be written as $\overline{z_1}, \overline{z_2}, \overline{z_3}, \overline{z_4}$. With this observation, it's easy to see that the polynomials $P(z)$ and $Q(4i\hspace{1pt}z)$ have the same roots. Hence, there exists some constant $K$ such that\begin{align*}P(z)=K*Q(4i\hspace{1pt}z)\end{align*}


 By comparing coefficients, its easy to see that $K=\frac{1}{(4i)^4}$. Hence $\frac{B*(4i)^2}{(4i)^4}=3$ and $\frac{D}{(4i)^4}=1$. Hence $B=-48$, $D=256$, so $B+D=208$ and our answer is $\boxed{(\textbf{D}) \: 208}$.

 


	\item (\textit{AMC 12 12A Fall 2021}) An organization has $30$ employees, $20$ of whom have a brand A computer while the other $10$ have a brand B computer. For security, the computers can only be connected to each other and only by cables. The cables can only connect a brand A computer to a brand B computer. Employees can communicate with each other if their computers are directly connected by a cable or by relaying messages through a series of connected computers. Initially, no computer is connected to any other. A technician arbitrarily selects one computer of each brand and installs a cable between them, provided there is not already a cable between that pair. The technician stops once every employee can communicate with each other. What is the maximum possible number of cables used?


 $\textbf{(A)}\ 190  \qquad\textbf{(B)}\  191 \qquad\textbf{(C)}\  192 \qquad\textbf{(D)}\  195 \qquad\textbf{(E)}\ 196$

 \textbf{Answer: }B



\textbf{Solution}

We claim that to maximize the number of cables used, we isolate one computer and connect all cables for the remaining $29$ computers, then connect one more cable for the isolated computer.

 If a brand A computer is isolated, then the technician can use at most $19\cdot10+1=191$ cables. If a brand B computer is isolated instead, then the technician can use at most $20\cdot9+1=181$ cables. Therefore, the answer is $\boxed{\textbf{(B)}\  191}.$
<u>\textbf{Remark}</u>

 Suppose that the computers are labeled $A_1, A_2, \ldots, A_{20}$ and $B_1, B_2, \ldots, B_{10}.$ We will prove the claim in the first paragraph of this solution:

 
            \begin{enumerate} 
        \item \textbf{With exactly $\boldsymbol{190}$ cables, it is possible that some computers cannot communicate with each other.}From this solution, it is clear that this claim is true: We isolate $A_{20}$ and connect all cables for $A_1, A_2, \ldots, A_{19}$ and $B_1, B_2, \ldots, B_{10}$ to get $19\cdot10=190$ cables. However, $A_{20}$ cannot communicate with any of these $29$ computers. 

  

 
        \item \textbf{With exactly $\boldsymbol{191}$ cables, all computers must be able to communicate with each other.}By the Pigeonhole Principle, we conclude that at least one brand B computer connects to all $20$ brand A computers. Without loss of generality, we assume that $B_1$ connects to all $20$ brand A computers. On the other hand, we conclude that at least $11$ brand A computers connect to all $10$ brand B computers. Without loss of generality, we assume that $A_1, A_2, \ldots, A_{11}$ connect to all $10$ brand B computers. 

 Therefore, any two computers must be able to communicate with each other:

 
            \begin{itemize} 
        \item Between $A_i$ and $A_j,$ where $i,j\in\{1,2,\ldots,20\}$ and $i\neq j$The relay is $A_i\leftrightarrow B_1\leftrightarrow A_j.$

        \item Between $B_i$ and $B_j,$ where $i,j\in\{1,2,\ldots,10\}$ and $i\neq j$The relay is $B_i\leftrightarrow \{A_1, A_2, \ldots, A_{11}\}\leftrightarrow B_j.$

        \item Between $A_i$ and $B_j,$ where $i\in\{1,2,\ldots,20\}$ and $j\in\{1,2,\ldots,10\}$The relay is $A_i\leftrightarrow B_1\leftrightarrow \{A_1, A_2, \ldots, A_{11}\}\leftrightarrow B_j.$

            \end{itemize}
            
            \end{enumerate}
            


	\item (\textit{AMC 12 12A Fall 2021}) For how many ordered pairs $(b,c)$ of positive integers does neither $x^2+bx+c=0$ nor $x^2+cx+b=0$ have two distinct real solutions?


 $\textbf{(A) } 4 \qquad \textbf{(B) } 6 \qquad \textbf{(C) } 8 \qquad \textbf{(D) } 12 \qquad \textbf{(E) } 16 \qquad$

 \textbf{Answer: }B



\textbf{Solution 1 (Casework)}

A quadratic equation does not have two distinct real solutions if and only if the \href{https://artofproblemsolving.com/wiki/index.php?title=Discriminant}{discriminant} is nonpositive. We conclude that:

 
            \begin{enumerate} 
        \item Since $x^2+bx+c=0$ does not have real solutions, we have $b^2\leq 4c.$

        \item Since $x^2+cx+b=0$ does not have real solutions, we have $c^2\leq 4b.$

            \end{enumerate}
            Squaring the first inequality, we get $b^4\leq 16c^2.$ Multiplying the second inequality by $16,$ we get $16c^2\leq 64b.$ Combining these results, we get 
 \[b^4\leq 16c^2\leq 64b.\]Since $b^4\leq 64b,$ it follows that $b\leq4.$ We apply casework to the value of $b:$

            \begin{itemize}
        \item If $b=1,$ then $1\leq 16c^2\leq 64,$ from which $c=1,2.$
            \end{itemize}
            
            \begin{itemize}
        \item If $b=2,$ then $16\leq 16c^2\leq 128,$ from which $c=1,2.$
            \end{itemize}
            
            \begin{itemize}
        \item If $b=3,$ then $81\leq 16c^2\leq 192,$ from which $c=3.$
            \end{itemize}
            
            \begin{itemize}
        \item If $b=4,$ then $256\leq 16c^2\leq 256,$ from which $c=4.$
            \end{itemize}
            Together, there are $\boxed{\textbf{(B) } 6}$ ordered pairs $(b,c),$ namely $(1,1),(1,2),(2,1),(2,2),(3,3),$ and $(4,4).$




\textbf{Solution 2 (Graphing)}

Similar to Solution 1, use the discriminant to get $b^2\leq 4c$ and $c^2\leq 4b$. These can be rearranged to $c\geq \frac{1}{4}b^2$ and $b\geq \frac{1}{4}c^2$. Now, we can roughly graph these two inequalities, letting one of them be the $x$ axis and the other be $y$. The graph of solutions should be above the parabola and under its inverse, meaning we want points on the graph or in the first area enclosed by the two graphs:
                    \begin{figure}[H]
                        \centering
                        \includegraphics[width=0.5\textwidth]{Solution/AMC_12_12A_Fall_2021_17_Solution_Figure_1.png}
                    \end{figure}
                    We are looking for lattice points (since $b$ and $c$ are positive integers), of which we can count $\boxed{\textbf{(B) } 6}$.

 



\textbf{Solution 3 (Graphing)}

We need to solve the following system of inequalities:
 \[ \left\{ \begin{array}{ll} b^2 - 4 c \leq 0 \\ c^2 - 4 b \leq 0 \end{array} \right.. \]Feasible solutions are in the region formed between two parabolas $b^2 - 4 c = 0$ and $c^2 - 4 b = 0$.

 Define $f \left( b \right) = \frac{b^2}{4}$ and $g \left( b \right) = 2 \sqrt{b}$.Therefore, all feasible solutions are in the region formed between the graphs of these two functions.

 For $b = 1$, we have $f(b) = \frac{1}{4}$ and $g(b) = 2$.Hence, the feasible $c$ are $1, 2$.

 For $b = 2$, we have $f(b) = 1$ and $g(b) = 2 \sqrt{2}$.Hence, the feasible $c$ are $1, 2$.

 For $b = 3$, we have $f(b) = \frac{9}{4}$ and $g(b) = 2 \sqrt{3}$.Hence, the feasible $c$ is $3$.

 For $b = 4$, we have $f(b) = 4$ and $g(b) = 4$.Hence, the feasible $c$ is $4$.

 For $b > 4$, we have $f(b) > g(b)$. Hence, there is no feasible $c$.

 Putting all cases together, the correct answer is $\boxed{\textbf{(B) } 6}$.

 



\textbf{Solution 4 (Oversimplified but Risky)}

A quadratic equation $Ax^2+Bx+C=0$ has one real solution if and only if $B^2-4AC=0.$ Similarly, it has imaginary solutions if and only if $B^2-4AC<0.$ We proceed as following:

 We want both $x^2+bx+c$ to be $1$ value or imaginary and $x^2+cx+b$ to be $1$ value or imaginary. $x^2+4x+4$ is one such case since $\sqrt {b^2-4ac}$ is $0.$ Also, $x^2+3x+3, x^2+2x+2, x^2+x+1$ are always imaginary for both $b$ and $c.$ We also have $x^2+x+2$ along with $x^2+2x+1$ since the latter has one solution, while the first one is imaginary. Therefore, we have $\boxed{\textbf{(B) } 6}$ total ordered pairs of integers.

 



\textbf{Solution 5 (Quick and Easy)}

We see that $b^2 \leq 4c$ and $c^2 \leq 4b.$ WLOG, assume that $b \geq c.$ Then we have that $b^2 \leq 4c \leq 4b$, so $b^2 \leq 4b$ and therefore $b \leq 4$, also meaning that $c \leq 4.$ This means that we only need to try 16 cases. Now we can get rid of the assumption that $b \geq c$, because we want ordered pairs. For $b = 1$ and $b = 2$, $c = 1$ and $c = 2$ work. When $b = 3$, $c$ can only be $3$, and when $b = 4$, only $c = 4$ works, for a total of $\boxed{\textbf{(B) } 6}$ ordered pairs of integers.

 



\textbf{Solution 6 (Fastest)}

We need both $b^2\leq 4c$ and $c^2\leq 4b$. 

 If $b=c$ then the above become $b^2\leq 4b\iff b\leq 4$, so we have four solutions $(k,k)$, where $k=1$,$2$,$3$,$4$.

 If $b<c$ then we only need $c^2\leq 4b$ since it implies $b^2< 4c$. Now$c^2\leq 4b\leq 4(c-1) \implies (c-2)^2\leq 0 \implies c=2$, so $b=1$. We plug $b=1$, $c=2$ back into $c^2\leq 4b$ and it works. So there is another solution $(1,2)$.

 By symmetry, if $b>c$ then $(b,c)=(2,1)$.

 Therefore the total number of solutions is $\boxed{\textbf{(B) } 6}$.

 



\textbf{Solution 7 (Shortest)}

Since $b^{2} - 4c \le 0$ and $c^{2} - 4b \le 0$, adding the two together yields $b^{2} + c^{2} \le 4(c+b)$. Obviously, this is not true if either $b$ or $c$ get too large, and they are equal when $b = c = 4$, so the greatest pair is $(4,4)$ and both numbers must be lesser for further pairs. For there to be two distinct real solutions, we can test all these pairs where $(b,c)$ are less than 4 (except for the already valid solution) on the original quadratics, and we find the working pairs are $(1,1)$, $(2,1)$, $(1,2)$, $(2,2)$, $(3,3)$, $(4,4)$ meaning there are $\boxed{\textbf{(B) } 6}$ pairs.

 - youtube.com/indianmathguy

 


	\item (\textit{AMC 12 12A Fall 2021}) Each of $20$ balls is tossed independently and at random into one of $5$ bins. Let $p$ be the probability that some bin ends up with $3$ balls, another with $5$ balls, and the other three with $4$ balls each. Let $q$ be the probability that every bin ends up with $4$ balls. What is $\frac{p}{q}$?


 $\textbf{(A)}\ 1 \qquad\textbf{(B)}\  4 \qquad\textbf{(C)}\  8 \qquad\textbf{(D)}\  12 \qquad\textbf{(E)}\ 16$

 \textbf{Answer: }E



\textbf{Solution 1 (Multinomial Coefficients)}

For simplicity purposes, we assume that the balls and the bins are both distinguishable.

 Recall that there are $5^{20}$ ways to distribute $20$ balls into $5$ bins. For $p,$ we choose one of the $5$ bins to have $3$ balls and another one of the $4$ bins to have $5$ balls. We get
 \[p=\frac{5\cdot4\cdot\binom{20}{3,5,4,4,4}}{5^{20}} \text{ and } q=\frac{\binom{20}{4,4,4,4,4}}{5^{20}}.\] Therefore, the answer is 
 \[\frac pq=\frac{5\cdot4\cdot\binom{20}{3,5,4,4,4}}{\binom{20}{4,4,4,4,4}}=\frac{5\cdot4\cdot\frac{20!}{3!\cdot5!\cdot4!\cdot4!\cdot4!}}{\frac{20!}{4!\cdot4!\cdot4!\cdot4!\cdot4!}}=\frac{5\cdot4\cdot(4!\cdot4!\cdot4!\cdot4!\cdot4!)}{3!\cdot5!\cdot4!\cdot4!\cdot4!}=\frac{5\cdot4\cdot4}{5}=\boxed{\textbf{(E)}\ 16}.\]~MRENTHUSIASM ~Jesshuang ~mathboy282

 



\textbf{Solution 2}

For simplicity purposes, we assume that the balls and the bins are both distinguishable.

 Let $q=\frac{x}{a},$ where $a$ is the total number of combinations and $x$ is the number of cases where every bin ends up with $4$ balls.

 We can take $1$ ball from one bin and place it in another bin so that some bin ends up with $3$ balls, another with $5$ balls, and the other three with $4$ balls each. Note that one configuration of $4{-}4{-}4{-}4{-}4$ corresponds to $5\cdot4\cdot4=80$ configurations of $3{-}5{-}4{-}4{-}4.$ On the other hand, one configuration of $3{-}5{-}4{-}4{-}4$ corresponds to $5$ configurations of $4{-}4{-}4{-}4{-}4.$

 Therefore, we have 
 \[p = \frac{80}{5}\cdot\frac{x}{a} = 16\cdot\frac{x}{a},\] from which $\frac{p}{q} = \boxed{\textbf{(E)}\ 16}.$




\textbf{Solution 3 (Binomial Coefficients)}

Since both of the cases will have $3$ bins with $4$ balls in them, we can leave those out. There are $2 \cdot \binom {5}{2} = 20$ ways to choose where to place the $3$ and the $5$. After that, there are $\binom {8}{3} = 56$ ways to put the $3$ and $5$ balls being put into the bins. For the $4,4,4,4,4$ case, after we canceled the $4,4,4$ out, we have $\binom {8}{4} = 70$ ways to put the $4$ balls inside the bins. Therefore, we have $\frac {56\cdot 20}{70}$ which is equal to $8 \cdot 2 = \boxed{\textbf{(E)}\ 16}$.

 



\textbf{Solution 4 (Set Theory / Graph Theory)}

Construct the set $A$ consisting of all possible $3{-}5{-}4{-}4{-}4$ bin configurations, and construct set $B$ consisting of all possible $4{-}4{-}4{-}4{-}4$ configurations. If we let $N$ be the total number of configurations possible, it's clear we want to solve for $\frac{p}{q} = \frac{\frac{|A|}{N}}{\frac{|B|}{N}} = \frac{|A|}{|B|}$. 

 Consider drawing an edge between an element in $A$ and an element in $B$ if it is possible to reach one configuration from the other by moving a single ball (Note this process is reversible.). Let us consider the total number of edges drawn.

 For any element in $A$, we may choose one of the $5$ balls in the $5$-bin and move it to the $3$-bin to get a valid element in $B$. This implies the number of edges is $5|A|$. 

 On the other hand, for any element in $B$, we may choose one of the $20$ balls and move it to one of the other $4$-bins to get a valid element in $A$. This implies the number of edges is $80|B|$.

 We equate the expressions to get $5|A| = 80|B|$, from which $\frac{|A|}{|B|} = \frac{80}{5} = \boxed{\textbf{(E)}\ 16}$.

 


	\item (\textit{AMC 12 12A Fall 2021}) Let $x$ be the least real number greater than $1$ such that $\sin(x) = \sin(x^2)$, where the arguments are in degrees. What is $x$ rounded up to the closest integer?


 $\textbf{(A) } 10 \qquad \textbf{(B) } 13 \qquad \textbf{(C) } 14 \qquad \textbf{(D) } 19 \qquad \textbf{(E) } 20$

 \textbf{Answer: }B



\textbf{Solution 1}

The smallest $x$ to make $\sin(x) = \sin(x^2)$ would require $x=x^2$, but since $x$ needs to be greater than $1$, these solutions are not valid.

 The next smallest $x$ would require $x=180-x^2$, or $x^2+x=180$. 

 After a bit of guessing and checking, we find that $12^2+12=156$, and $13^2+13=182$, so the solution lies between $12{ }$ and $13$, making our answer $\boxed{\textbf{(B) } 13}.$

 Note: One can also solve the quadratic and estimate the radical.

 Note: the answer is 13 because 182 is far more close to 180 than 156 and as x is positive then it is much closer, so the answer is 13

 Note: solving: x\textasciicircum{}2+x-180=0 so x= (-1+- sqrt(721))/2 clearly only x= (-1+sqrt(721))/2 is feasible here. we make a rough estimate of 721 being 729 (as 729 is the closet perfect square),  so it is (-1+27)/2=13.    in addition, we could try to estimate with 12 (which is the worth case),   26\textasciicircum{}2=4*169 =676 which yeilds us x=(-1+26)/2=12.5 which closely, it is still 13 (12.5 is close to 13) 

 therefore the answer is undeinablely 13. majorly it is some addition onto the original basis

 



\textbf{Solution 2}

For choice $\textbf{(A)},$ we have
 \begin{align*} \left| \sin x - \sin \left( x^2 \right) \right| & = \left| \sin 10^\circ - \sin \left( \left( 10^2 \right)^\circ \right) \right| \\ & = \left| \sin 10^\circ - \sin 100^\circ \right| \\ & = \left| \sin 10^\circ - \sin 80^\circ \right| \\ & = \sin 80^\circ - \sin 10^\circ . \end{align*}
For choice $\textbf{(B)},$ we have
 \begin{align*} \left| \sin x - \sin \left( x^2 \right) \right| & = \left| \sin 13^\circ - \sin \left( \left( 13^2 \right)^\circ \right) \right| \\ & = \left| \sin 13^\circ - \sin 169^\circ \right| \\ & = \left| \sin 10^\circ - \sin 11^\circ \right| \\ & = \sin 11^\circ - \sin 10^\circ . \end{align*}
For choice $\textbf{(C)},$ we have
 \begin{align*} \left| \sin x - \sin \left( x^2 \right) \right| & = \left| \sin 14^\circ - \sin \left( \left( 14^2 \right)^\circ \right) \right| \\ & = \left| \sin 14^\circ - \sin 196^\circ \right| \\ & = \left| \sin 14^\circ + \sin 16^\circ \right| \\ & = \sin 14^\circ + \sin 16^\circ . \end{align*}
For choice $\textbf{(D)},$ we have
 \begin{align*} \left| \sin x - \sin \left( x^2 \right) \right| & = \left| \sin 19^\circ - \sin \left( \left( 19^2 \right)^\circ \right) \right| \\ & = \left| \sin 19^\circ - \sin 361^\circ \right| \\ & = \left| \sin 19^\circ - \sin 1^\circ \right| \\ & = \sin 19^\circ - \sin 1^\circ . \end{align*}
For choice $\textbf{(E)},$ we have
 \begin{align*} \left| \sin x - \sin \left( x^2 \right) \right| & = \left| \sin 20^\circ - \sin \left( \left( 20^2 \right)^\circ \right) \right| \\ & = \left| \sin 20^\circ - \sin 400^\circ \right| \\ & = \left| \sin 20^\circ - \sin 40^\circ \right| \\ & = \sin 40^\circ - \sin 20^\circ . \end{align*}
Therefore, the answer is $\boxed{\textbf{(B) }13},$ as $\sin 11^\circ - \sin 10^\circ$ is the closest to $0.$

 \textbf{Easy Solution with Intermediate Value Theorem} We know that $x^2-x=180-2x+360k$ or $x^2-x=360k$ from looking at the $\text{sin}$ values on the unit circle. With these equations, we can just insert the values that are given in the $5$ answers. 

 
Let's try $\text{A}$ first. 

 We know that one of the equations should be greater on the left, and another equation we make should be greater on the right for the solution to be in between the values inserted into the equations. The two numbers it is in between is $9$ and $10$ since it is rounded up. 

 
$100-10$ vs $180-20,$
$81-9$ vs $180-18.$

 
This does not satisfy what we would like to find.

 
Now let's try $\text{B}$.

 Between $12$ and $13$;

 
$144-12$ vs $180-24,$
$169-13$ vs $180-26.$

 
Now this satisfies our hunt for the solution.The answer is $\boxed{\text{(B)}}$

 emilyyunhanq@gmail.com

 Solution by Emily Q

 
(edited by A7456321)

 


	\item (\textit{AMC 12 12A Fall 2021}) For each positive integer $n$, let $f_1(n)$ be twice the number of positive integer divisors of $n$, and for $j \ge 2$, let $f_j(n) = f_1(f_{j-1}(n))$. For how many values of $n \le 50$ is $f_{50}(n) = 12?$
$\textbf{(A) }7\qquad\textbf{(B) }8\qquad\textbf{(C) }9\qquad\textbf{(D) }10\qquad\textbf{(E) }11$

 \textbf{Answer: }D



\textbf{Solution 1}

First, we can test values that would make $f(x)=12$ true. For this to happen $x$ must have $6$ divisors, which means its prime factorization is in the form $pq^2$ or $p^5$, where $p$ and $q$ are prime numbers. Listing out values less than $50$ which have these prime factorizations, we find $12,18,20,28,44,45,50$ for $pq^2$, and just $32$ for $p^5$. Here $12$ especially catches our eyes, as this means if one of $f_i(n)=12$, each of $f_{i+1}(n), f_{i+2}(n), ...$ will all be $12$. This is because $f_{i+1}(n)=f(f_i(n))$ (as given in the problem statement), so were $f_i(n)=12$, plugging this in we get $f_{i+1}(n)=f(12)=12$, and thus the pattern repeats. Hence, as long as for a $i$, such that $i\leq 50$ and $f_{i}(n)=12$, $f_{50}(n)=12$ must be true, which also immediately makes all our previously listed numbers, where $f_1(x)=12$, possible values of $n$.

 We also know that if $f_1(x)$ were to be any of these numbers, $x$ would satisfy $f_{50}(n)$ as well. Looking through each of the possibilities aside from $12$, we see that $f_1(x)$ could only possibly be equal to $20$ and $18$, and still have $x$ less than or equal to $50$. This is because if $f_1(x)=20$ or $18$, $f_2 (x) = 12$. This would mean $x$ must have $10$, or $9$ divisors, and testing out, we see that $x$ will then be of the form $p^4q$, or $p^2q^2$. The only two values less than or equal to $50$ would be $48$ and $36$ respectively. From here there are no more possible values, so tallying our possibilities we count $\boxed{\textbf{(D) }10}$ values (Namely $12,18,20,28,32,36,44,45,48,50$).

 



\textbf{Solution 2 (Rigorous reasoning on why there cannot be any other solutions)}

First, take note that the maximum possible value of $f_1(n)$ for $1 \le n \le k$ increases as $k$ increases. When we increase $k$, we either include a new value of $n$ that has more divisors than any previous values of $n$, increasing the maximum of $f_1(x)$, or none of the new $n$ have more divisors, which would leave the maximum of $f_1(x)$ the same as before. In other words, $f_1(x)$ is increasing. This means that $f_1(x)$ decreases when we decrease $k$. 

 Also, let $f_1(n) = 2d(n)$ where $d(n)$ is the number of divisors of n.

 Since $n \le 50$, $f_1(n) \le 20$. This maximum occurs when $d(n) = 10 \implies n = 2^4 \cdot 3 = 48$. Next, since $f_1(n) \le 20$, $f_1(f_1(n)) \le 12 \implies f_2(n) \le 12$. This maximum occurs when $d(f_1(n)) = 6 \implies n = 2 \cdot 3^2 = 18, n = 2^2 \cdot 3 = 12$. Since $f_2(n) \le 12$, $f_1(f_2(n)) \le 12 \implies f_3(n) \le 12$, once again. This maximum again occurs when $d(f_2(n)) = 6 \implies f_2(n) = 2^2 \cdot 3 = 12$. Notice that, from this point on, for $k \geq 3$, $f_k(n) \leq 12$ by a straightforward induction argument.

 On a hunch, we investigate what happens when $f_2(n) < 12$, as this is the first time when $f_k(n)$ is bounded to a maximum of 12.Suppose for the sake of contradiction that $f_2(n) < 12$. Then, $f_3(n) < 12$, as $f_2(n) = 12$ is the only number that would maximize $f_3(n))$ for $f_2(n) \le 12$. As a result, since $f_1(n)$ is increasing, and because $12$ is the only value for which $f_1$  is maximized, $f_1$ "steps down" from a maximum of $6 \cdot 2 = 12$ if $n < 12$. This means we must have that $f_1(f_3(n)) < f_1(12) = 12 \implies f_4(n) < 12$. We continue applying $f_1$ on both sides (which is possible since $f_1$ is increasing) until we reach $f_{50}$, giving us that $f_{50}(n) < 12$. However, $f_{50}(n) = 12$, which is a contradiction. Thus, $f_2(n) = 12$. 

 Now, we can finally solve for the solutions with a clean conscience - we proceed similarly to Solution 1. $f_2(n) = 12 \implies f_1(f_1(n)) = 12 \implies d(f_1(n)) = 6$. $d(f_1(n)) = 6$. This gives us two cases. First, we have the case where $f_1(n) = p^2 \cdot q$ where $p$ and $q$ are primes. Second, we have the case where $f_1(n)=p^5$ where $p$ is a prime.  For both cases, since $f_1(n) \le 20$, $f_1(n)$ can only be $12$, $18$, or $20$.

 If $f_1(n) = 12$, then $d(n) = 6 \implies n = p^5, p^2 \cdot q \implies n \in \{ 12, 18, 20, 28, 32, 44, 45, 50 \}$, resulting in 8 solutions. 

 If $f_1(n) = 18$, then $d(n) = 9 \implies n = p^8, p^2 \cdot q^2 \implies n = 36$, giving us one more solution. Finally, $f_1(n) = 20 \implies d(n) = 10 \implies n = p^9, p^4 \cdot q \implies n = 48$. 

 Thus, in total, we have $\boxed{\textbf{(D)} 10}$ solutions.

 



\textbf{Solution 3}

$\textbf{Observation 1}$: $f_1 \left( 12 \right) = 12$.

 Hence, if $n$ has the property that $f_j \left( n \right) = 12$ for some $j$, then $f_k \left( n \right) = 12$ for all $k > j$.

 $\textbf{Observation 2}$: $f_1 \left( 8 \right) = 8$.

 Hence, if $n$ has the property that $f_j \left( n \right) = 8$ for some $j$, then $f_k \left( n \right) = 8$ for all $k > j$.

 $\textbf{Case 1}$: $n = 1$.

 We have $f_1 \left( n \right) = 2$, $f_2 \left( n \right) = f_1 \left( 2 \right) = 4$, $f_3 \left( n \right) = f_1 \left( 4 \right) = 6$, $f_4 \left( n \right) = f_1 \left( 6 \right) = 8$. Hence, Observation 2 implies $f_{50} \left( n \right) = 8$.

 $\textbf{Case 2}$: $n$ is prime.

 We have $f_1 \left( n \right) = 4$, $f_2 \left( n \right) = f_1 \left( 4 \right) = 6$, $f_3 \left( n \right) = f_1 \left( 6 \right) = 8$. Hence, Observation 2 implies $f_{50} \left( n \right) = 8$.

 $\textbf{Case 3}$: The prime factorization of $n$ takes the form $p_1^2$.

 We have $f_1 \left( n \right) = 6$, $f_2 \left( n \right) = f_1 \left( 6 \right) = 8$. Hence, Observation 2 implies $f_{50} \left( n \right) = 8$.

 $\textbf{Case 4}$: The prime factorization of $n$ takes the form $p_1^3$.

 We have $f_1 \left( n \right) = 8$. Hence, Observation 2 implies $f_{50} \left( n \right) = 8$.

 $\textbf{Case 5}$: The prime factorization of $n$ takes the form $p_1^4$.

 We have $f_1 \left( n \right) = 10$, $f_2 \left( n \right) = f_1 \left( 10 \right) = 8$. Hence, Observation 2 implies $f_{50} \left( n \right) = 8$.

 $\textbf{Case 6}$: The prime factorization of $n$ takes the form $p_1^5$.

 We have $f_1 \left( n \right) = 12$. Hence, Observation 1 implies $f_{50} \left( n \right) = 12$.

 In this case the only $n$ is $2^5 = 32$.

 $\textbf{Case 7}$: The prime factorization of $n$ takes the form $p_1 p_2$.

 We have $f_1 \left( n \right) = 8$. Hence, Observation 2 implies $f_{50} \left( n \right) = 8$.

 $\textbf{Case 8}$: The prime factorization of $n$ takes the form $p_1 p_2^2$.

 We have $f_1 \left( n \right) = 12$. Hence, Observation 1 implies $f_{50} \left( n \right) = 12$.

 In this case, all $n$ are $12, 18, 20, 28, 44, 45,$ and $50$.

 $\textbf{Case 9}$: The prime factorization of $n$ takes the form $p_1 p_2^3$.

 We have $f_1 \left( n \right) = 16$, $f_2 \left( n \right) = f_1 \left( 16 \right) = 10$, $f_3 \left( n \right) = f_1 \left( 10 \right) = 8$. Hence, Observation 2 implies $f_{50} \left( n \right) = 8$.

 $\textbf{Case 10}$: The prime factorization of $n$ takes the form $p_1 p_2^4$.

 We have $f_1 \left( n \right) = 20$, $f_2 \left( n \right) = f_1 \left( 20 \right) = 12$. Hence, Observation 1 implies $f_{50} \left( n \right) = 12$.

 In this case, the only $n$ is $48$.

 $\textbf{Case 11}$: The prime factorization of $n$ takes the form $p_1^2 p_2^2$.

 We have $f_1 \left( n \right) = 18$, $f_2 \left( n \right) = f_1 \left( 18 \right) = 12$. Hence, Observation 1 implies $f_{50} \left( n \right) = 12$.

 In this case, the only $n$ is $36$.

 $\textbf{Case 12}$: The prime factorization of $n$ takes the form $p_1 p_2 p_3$.

 We have $f_1 \left( n \right) = 16$, $f_2 \left( n \right) = f_1 \left( 16 \right) = 10$, $f_3 \left( n \right) = f_2 \left( 10 \right) = 8$. Hence, Observation 2 implies $f_{50} \left( n \right) = 8$.

 Putting all cases together, the number of feasible $n \leq 50$ is $\boxed{\textbf{(D) }10}$.

 


 


	\item (\textit{AMC 12 12A Fall 2021}) Let $ABCD$ be an isosceles trapezoid with $\overline{BC} \parallel \overline{AD}$ and $AB=CD$. Points $X$ and $Y$ lie on diagonal $\overline{AC}$ with $X$ between $A$ and $Y$, as shown in the figure. Suppose $\angle AXD = \angle BYC = 90^\circ$, $AX = 3$, $XY = 1$, and $YC = 2$. What is the area of $ABCD$?


 
                    \begin{figure}[H]
                        \centering
                        \includegraphics[width=0.5\textwidth]{Problem/AMC_12_12A_Fall_2021_21_Problem_Figure_1.png}
                    \end{figure}
                    
$\textbf{(A)}\: 15\qquad\textbf{(B)} \: 5\sqrt{11}\qquad\textbf{(C)} \: 3\sqrt{35}\qquad\textbf{(D)} \: 18\qquad\textbf{(E)} \: 7\sqrt{7}$

 \textbf{Answer: }C



\textbf{Solution 1}

First realize that $\triangle BCY \sim \triangle DAX.$ Thus, because $CY: XA = 2:3,$ we can say that $BY = 2s$ and $DX = 3s.$ From the Pythagorean Theorem, we have $AB^2 =(2s)^2 + 4^2 = 4s^2 + 16$ and $CD^2 = (3s)^2 + 3^2 = 9s^2 + 9.$ Because $AB = CD,$ from the problem statement, we have that 
 \[4s^2 + 16 = 9s^2 + 9.\] Solving, gives $s = \frac{\sqrt{7}}{\sqrt{5}}.$ To find the area of the trapezoid, we can compute the area of $\triangle ABC$ and add it to the area of $\triangle ACD.$ Thus, the area of the trapezoid is $\frac{1}{2} \cdot 2 \cdot \frac{\sqrt{7}}{\sqrt{5}} \cdot 6 + \left(\frac{1}{2} \cdot 3 \cdot \frac{\sqrt{7}}{\sqrt{5}} \cdot 6 \right) = \frac{15\sqrt{7}}{\sqrt{5}} = 3\sqrt{35}.$ Thus, the answer is $\boxed{\textbf{(C)} \: 3\sqrt{35}}.$




\textbf{Solution 2}

We put the graph to a coordinate plane. We put $X$ to the origin, $AC$ to the $x$-axis, and $DX$ to the $y$-axis.

 We have the coordinates of the following points: $X = \left( 0 , 0 \right)$, $A = \left( 3 , 0 \right)$, $Y = \left( - 1 , 0 \right)$, $C = \left( - 3 , 0 \right)$.

 Denote $BY = u$, $DX = v$. Hence, $B = \left( - 1 , u \right)$, $D = \left( 0 , - v \right)$.

 Hence, the slope of $BC$ is $m_{BC} = \frac{u}{2}$.The slope of $AD$ is $m_{AD} = \frac{v}{3}$.

 Because $BC \parallel AD$, $m_{BC} = m_{AD}$. Hence,
 \[ \frac{u}{2} = \frac{v}{3} . \hspace{1cm} (1) \]

 In $\triangle AYB$, following from the Pythagorean theorem, we have $AB^2 = AY^2 + BY^2 = 4^2 + u^2$.

 In $\triangle CXD$, following from the Pythagorean theorem, we have $CD^2 = CX^2 + XD^2 = 3^2 + v^2$.

 Because $AB = CD$,
 \[ 4^2 + u^2 = 3^2 + v^2 . \hspace{1cm} (2) \]

 Solving Equations (1) and (2), we get $u = \frac{2 \sqrt{35}}{5}$, $V = \frac{3 \sqrt{35}}{5}$.

 Therefore,
 \begin{align*} {\textrm Area} \ ABCD & = {\textrm Area} \ \triangle ABC + {\textrm Area} \ \triangle  ADC \\ & = \frac{1}{2} AC \cdot BY + \frac{1}{2} AC \cdot DX \\ & = \frac{1}{2} AC \left( BY + DY \right) \\ & = 3 \sqrt{35} . \end{align*}


 Therefore, the answer is $\boxed{\textbf{(C)} \: 3\sqrt{35}}$.

 



\textbf{Solution 3}

Since both $\overline{BY}$ and $\overline{DX}$ are perpendicular to $\overline{AC}$, we have $BD^2=XY^2+(BY+DX)^2$.

 Since the trapezoid is equilateral, $BD=AC=6$. Hence, $BY+DX=\sqrt{6^2-1^2}=\sqrt{35}$.

 Notice that both triangles $ABC$ and $ADC$ share a common base $\overline{AC}$.

 Therefore, 
 \[[ABCD]=[ABC]+[ADC]=\frac{1}{2}AC\cdot(BY+DX)=\frac{1}{2}\cdot6\cdot\sqrt{35}=\boxed{\textbf{(C)} \: 3\sqrt{35}}.\]~RunyangWang

 


	\item (\textit{AMC 12 12A Fall 2021}) Azar and Carl play a game of tic-tac-toe. Azar places an $X$ in one of the boxes in a $3$-by-$3$ array of boxes, then Carl places an $O$ in one of the remaining boxes. After that, Azar places an $X$ in one of the remaining boxes, and so on until all 9 boxes are filled or one of the players has 3 of their symbols in a row - horizontal, vertical, or diagonal - whichever comes first, in which case that player wins the game. Suppose the players make their moves at random, rather than trying to follow a rational strategy, and that Carl wins the game when he places his third $O$. How many ways can the board look after the game is over?


 $\textbf{(A) } 36 \qquad\textbf{(B) } 112 \qquad\textbf{(C) } 120 \qquad\textbf{(D) } 148 \qquad\textbf{(E) } 160$

 \textbf{Answer: }D



\textbf{Solution 1}

We need to find out the number of configurations with 3 $O$ and 3 $X$ with 3 $O$ in a row, and 3 $X$ not in a row.

 $\textbf{Case 1}$: 3 $O$ are in a horizontal row or a vertical row.

 Step 1: We determine the row that 3 $O$ occupy.

 The number of ways is 6.

 Step 2: We determine the configuration of 3 $X$.

 The number of ways is $\binom{6}{3} - 2 = 18$. (6 open spots to put the $X$ and subtracting cases where they fill row/column)

 In this case, following from the rule of product, the number of ways is $6 \cdot 18 = 108$.

 $\textbf{Case 2}$: 3 $O$ are in a diagonal row.

 Step 1: We determine the row that 3 $O$ occupy.

 The number of ways is 2.

 Step 2: We determine the configuration of 3 $X$.

 The number of ways is $\binom{6}{3}  = 20$.

 In this case, following from the rule of product, the number of ways is $2 \cdot 20 = 40$.

 Putting all cases together, the total number of ways is $108 + 40 = 148$.

 Therefore, the answer is $\boxed{\textbf{(D) }148}$.

 



\textbf{Solution 2 (Complementary)}

First count all the ways to have 3 $O$ in a row and 3 $X$. The $O$ can be arranged in 8 possible ways ($3$ horizontal, $3$ vertical, $2$ diagonal), and in each case, the $X$ can be arranged in $\binom{6}{3}$ ways, so there are $8* \binom{6}{3}$ = 160 arrangements.

 Now exclude the ones in which the $X$ are also in a row column. Either the $X$ and $O$ both form rows, which can be achieved in 6 ways, or they both form columns, which can also be acheived in 6 ways. Excluding these 12 arrangements, the answer is shown to be $\boxed{\textbf{(D) }148}$.

 


	\item (\textit{AMC 12 12A Fall 2021}) A quadratic polynomial with real coefficients and leading coefficient $1$ is called $\emph{disrespectful}$ if the equation $p(p(x))=0$ is satisfied by exactly three real numbers. Among all the disrespectful quadratic polynomials, there is a unique such polynomial $\tilde{p}(x)$ for which the sum of the roots is maximized. What is $\tilde{p}(1)$?


 $\textbf{(A) } \frac{5}{16} \qquad\textbf{(B) } \frac{1}{2} \qquad\textbf{(C) } \frac{5}{8} \qquad\textbf{(D) } 1 \qquad\textbf{(E) } \frac{9}{8}$

 \textbf{Answer: }A



\textbf{Solution 1 (Standard Form)}

Let $r_1$ and $r_2$ be the roots of $\tilde{p}(x)$. Then, $\tilde{p}(x)=(x-r_1)(x-r_2)=x^2-(r_1+r_2)x+r_1r_2$. The solutions to $\tilde{p}(\tilde{p}(x))=0$ is the union of the solutions to 
 \[\tilde{p}(x)-r_1=x^2-(r_1+r_2)x+(r_1r_2-r_1)=0\] and 
 \[\tilde{p}(x)-r_2=x^2-(r_1+r_2)x+(r_1r_2-r_2)=0.\] Note that one of these two quadratics has one solution (a double root) and the other has two as there are exactly three solutions. WLOG, assume that the quadratic with one root is $x^2-(r_1+r_2)x+(r_1r_2-r_1)=0$. Then, the discriminant is $0$, so $(r_1+r_2)^2 = 4r_1r_2 - 4r_1$. Thus, $r_1-r_2=\pm 2\sqrt{-r_1}$, but for $x^2-(r_1+r_2)x+(r_1r_2-r_2)=0$ to have two solutions, it must be the case that $r_1-r_2=- 2\sqrt{-r_1} (*)$. It follows that the sum of the roots of $\tilde{p}(x)$ is $2r_1 + 2\sqrt{-r_1}$, whose maximum value occurs when $r_1 = - \frac{1}{4} (\star)$. Solving for $r_2$ yields $r_2 = \frac{3}{4}$. Therefore, $\tilde{p}(x)=x^2 - \frac{1}{2} x - \frac{3}{16}$, so $\tilde{p}(1)= \boxed{\textbf{(A) } \frac{5}{16}}$.

 <u>\textbf{Remarks}</u>

            \begin{itemize}
        \item For $x^2-(r_1+r_2)x+(r_1r_2-r_2)=0$ to have two solutions, the discriminant $(r_1+r_2)^2-4r_1r_2+4r_2$ must be positive. From here, we get that $(r_1-r_2)^2>-4r_2$, so $-4r_1>-4r_2 \implies r_1<r_2$. Hence, $r_1-r_2$ is negative, so $r_1-r_2=-2\sqrt{-r_1}$.
            \end{itemize}
            
            \begin{itemize}
        \item Set $\sqrt{-r_1}=x$. Now $r_1+\sqrt{-r_1}=-x^2+x$, for which the maximum occurs when $x=\frac{1}{2} \rightarrow r_1=-\frac{1}{4}$.
            \end{itemize}
            



\textbf{Solution 2 (Vertex Form)}

Let $p(x)=(x-h)^2+k$ for some real constants $h$ and $k.$ Suppose that $p(x)$ has real roots $r$ and $s.$

 Since $p(p(x))=0,$ we conclude that $p(x)=r$ or $p(x)=s.$ Without loss of generality, we assume that $p(x)=r$ has two real solutions and $p(x)=s$ has one real solution. Therefore, we have $k=s,$ from which 
 \[p(x)=(x-h)^2+s.\]As $p(s)=0,$ we expand the left side to obtain $(s-h)^2+s=0,$ or 
 \[s^2-(2h-1)s+h^2=0. \hspace{15mm}(\bigstar)\]Since $(\bigstar)$ has real solutions for $s,$ the discriminant is nonnegative: $(2h-1)^2-4h^2\geq0.$ We solve this inequality to get $h\leq \frac14.$

 Either by the axis of symmetry or Vieta's formulas, note that $r+s=2h.$ As we wish to maximize $2h,$ we maximize $h.$ Substituting $h=\frac14$ into $(\bigstar),$ we obtain $s^2+\frac12s+\frac{1}{16}=0.$ We factor the left side to get $\left(s+\frac14\right)^2=0,$ or $s=-\frac14.$

 Finally, the unique such polynomial is 
 \[\tilde{p}(x)=\left(x-\frac14\right)^2-\frac14,\] from which $\tilde{p}(1)=\boxed{\textbf{(A) } \frac{5}{16}}.$




\textbf{Solution 2.5 (Vertex Form but Done Differently)}

Continuing from the previous solution's reasoning, we know that the $y$-coordinate of the vertex must also be a root of the polynomial. Thus, $\tilde{p}(x) = \left(x - \frac{s + r}{2}\right)^2 + s$, where $r$ and $s$ are roots of the polynomial.

 We also know that in order for $\tilde{p}(x) = r$ to have two solutions, $y = r$ must intersect $\tilde{p}(x)$ above its minimum/vertex. Thus, $r > s$. Let $r = s + a$. Then, $\tilde{p}(x) = \left(x - \frac {2s + a}{2}\right)^2 + s$. 

 Since $s$ is a root, $\tilde{p}(s) = 0$ so $\left(s - \frac{2s + a}{2}\right)^2 +s = 0 \implies \frac{a^2}{4} + s = 0 \implies s = - \frac{1}{4} a^2$. We wish to maximize $s + s + a = 2s + a = - \frac{1}{2} a^2 + a$. 

 This quadratic is maximized at $a = 1$ (when converting the quadratic to vertex form). Then, since $\tilde{p}(s + a) = 0 \implies \tilde{p}(s+1) = 0$, $\left(s + 1 - s - \frac12\right)^2 + s = 0 \implies \frac{1}{4} + s = 0 \implies s = -\frac{1}{4}$. 

 Thus, $\tilde{p}(x) = \left(x - \frac{1}{4}\right)^2 - \frac{1}{4}$. So, $\tilde{p}(1) = \boxed{\textbf{(A) } \frac{5}{16}}$.

 



\textbf{Solution 3 (Symmetry)}

Let $\tilde p(x)=(x-h)^2+k$. We seek to maximize the sum of the roots $2h$, so we maximize $h$.

 Let $P(x)=\tilde p(\tilde p(x))=((x-h)^2+k-h)^2+k$. Note that $P(x)$ is symmetric about $x=h$. 

 $P(x)=0$ has $3$ real solutions, Due to the complex conjugate theorem, $P(x$) must have $4$ real roots. Therefore, $P(x)$ must have exactly $1$ double root. 

 This root cannot be to the left or to the right of $x=h$, as the symmetry of the function would mean that there would be another double root reflected across the $x=h$. It follows that the double root could only be situated at $x=h$. 

 $0=P(h)=((h-h)^2+k-h)^2+k=(k-h)^2+k$. Expanding and writing this out in terms of k, $k^2+(1-2h)k+h^2=0$. 

 In order for this to have a solution, the discriminant has to be non-negative. In other words, $(1-2h)^2-4h^2\geq0$. 

 This simplifies to $1-4h\geq0$, or $h\leq\frac{1}{4}$. 

 As we seek to maximize $h$, we set $h=\frac{1}{4}$ and see that $k=-\frac{1}{4}$. 

 Therefore, $\tilde p(x)=\left(x-\frac{1}{4}\right)^2-\frac{1}{4}$, and $\tilde p(1)=\left(1-\frac{1}{4}\right)^2-\frac{1}{4}=\frac{9}{16}-\frac{1}{4}=\boxed{\textbf{(A) } \frac{5}{16}}$.

 



\textbf{Solution 4 (Discriminant)}

Equation $p \left( p \left( x \right) \right) = 0$ is equivalent to the following system of equations:
 \begin{align*} p \left( u \right) &= 0, \\ p \left( x \right) - u &= 0. \end{align*}


 Denote $p \left( x \right) = x^2 + p x + q$.

 Denote by $r_1$ and $r_2$ two roots of $p \left( x \right) = 0$.

 Because $p \left( p \left( x \right) \right) = 0$ has three real roots, we must have the properties that $r_1$ and $_2$ are real with $r_1 \neq r_2$.Without loss of generality, we assume $r_1 < r_2$.

 We notice that all roots of $p \left( p \left( x \right) \right) = 0$are the collection of all roots of $p \left( x \right) - r_1 = 0$ and all roots of $p \left( x \right) - r_2 = 0$.

 Because each of these two equations is quadratic, it has two roots (may be identical).To get a total number of three roots, one equation must have two identical roots.

 Because $r_1 < r_2$, equation $p \left( x \right) - r_1 = 0$  has two identical roots.Hence, the discriminant of this equation satisfies
 \[ p^2 - 4 \left( q - r_1 \right) = 0 . \hspace{1cm} \]

 Because $r_1 < r_2$, $r_1 = \frac{- p - \sqrt{p^2 - 4 q}}{2}$.

 Hence, the above equation can be written as
 \[ p^2 - 4 \left( q - \frac{- p - \sqrt{p^2 - 4 q}}{2} \right) = 0 . \]

 This can be reorganized as
 \[ \left( p^2 - 4 q \right) - 2 \sqrt{p^2 - 4 q} - 2 p = 0 . \hspace{1cm} (1) \]

 Define $x = \sqrt{p^2 - 4 q}$.Hence, the value of $p$ should ensure that equation $x^2 - 2 x - 2 p = 0$ has at least one real nonnegative root.

 This condition can be satisfied if the discriminant of this equation is nonnegative: $\left( -2 \right)^2 - 4 \cdot 1 \cdot \left( - 2 p \right) \geq 0$.Hence, $p \geq - \frac{1}{2}$.

 Now, we are ready to find $\tilde p \left( x \right)$.

 Following from Vieta's formula, $r_1 + r_2 = - p$.Hence, to get $r_1 + r_2$ maximized, we need to find smallest $p$.

 Because $p \geq - \frac{1}{2}$, the smallest $p$ is $-\frac{1}{2}$.Plugging this value into Equation (1), we get $\sqrt{p^2 - 4 q} = 1$.Hence, $q = - \frac{3}{16}$.

 Thus, $\tilde p \left( x \right) = x^2 - \frac{1}{2} x - \frac{3}{16}$.Therefore, $\tilde p \left( 1 \right) = \boxed{\textbf{(A) } \frac{5}{16}}$.

 



\textbf{Solution 5 (Factored Form)}

The disrespectful function $p(x)$ has leading coefficient $1$, so it can be written in factored form as $(x-r)(x-s)$. Now the problem states that all $p(x)$ must satisfy $p(p(x)) = 0$. Plugging our form in, we get 
 \[((x-r)(x-s)-r)((x-r)(x-s)-s) = 0.\] The roots of this equation are $(x-r)(x-s) = r$ and $(x-r)(x-s) = s$. By the fundamental theorem of algebra, each root must have two roots for a total of four possible values of $x$ yet the problem states that this equation is satisfied by three values of $x$. Therefore one equation must give a double root. Without loss of generality, let the equation $(x-r)(x-s) = r$ be the equation that produces the double root. Expanding gives $x^2-(r+s)x+rs-r = 0$. We know that if there is a double root to this equation, the discriminant must be equal to zero, so $(r+s)^2-4(rs-r) = 0 \implies r^2+2rs+s^2-4rs+4r = 0 \implies r^2-2rs+s^2+4r = 0$.

 From here two solutions can progress.

 



\textbf{Solution 5.1 (Quadratic Formula)}

We can rewrite $r^2-2rs+s^2+4r = 0$ as $(r-s)^2+4r = 0$. Let's keep our eyes on the ball; we want to find the disrespectful quadratic that maximizes the sum of the roots, which is $r+s$. Let this be equal to a new variable, $m$, so that our problem is reduced to maximizing this variable. We can rewrite our equation in terms of $m$ as $(2r-m)^2 + 4r = 0 \implies m^2- 4rm + 4r^2+4r = 0$. 

 This is a quadratic in $m$, so we can use the quadratic formula: 
 \[m = \frac{4r \pm \sqrt{16r^2-4(4r^2+4r)}}{2} = 2r \pm \sqrt{-4r} = 2\left(r \pm \sqrt{-r}\right).\]It will be easier to think without the square root, so let $q = \sqrt{-r}$. We can rewrite the equation as $m = 2(-q^2 \pm q)$. We want to maximize $m$, so we take the plus value of the right-hand-side of the equation. Then,
 \[m=2(-q^2+q) \implies m = -2q(q-1).\]To maximize $m$, we find the vertex of the right-hand side of the equation. The vertex of $-2q(q-1)$ is the average of the roots of the equation which is $\frac{0+1}{2} = \frac{1}{2}$. This means that $r = -q^2 \implies \boldsymbol{r = -\frac{1}{4}}$ and $m = -2q(q-1) \implies m = \frac{1}{2}$. Therefore, $m-r = s \implies \boldsymbol{s = \frac{3}{4}}$.

 



\textbf{Solution 5.2 (Derivation-Rotated Conics)}

We see that the equation $r^2-2rs+s^2+4r = 0$ is in the form of the general equation of a rotated conic 
 \[Ax^2+Bxy+Cy^2+Dx+Ey+F = 0.\] Because $B^2 -4AC = (-2)^2 - 4(1)(1) = 0$, this rotated conic is a parabola. 

 The definition of a parabola is the locus of all points that are equidistant from a point (focus) and line (directrix). Let the focus and directrix of this particular parabola be $(a,b)$ and $y = mx+c$. Then we can try to find the general form of a rotated parabola in terms of $a, b, m,$and$c$. 

 The distance between two points $(x,y)$ and $(a,b)$ is $\sqrt{(x-a)^2+(y-b)^2}$. Therefore this is the distance from any point on the parabola to the focus. 

 The distance from a point $(x,y)$ to a line $y = mx+c \implies mx-y+c = 0$ is $\frac{|mx-y+c|}{\sqrt{m^2+1}}$.

 We can set these two equal to each other and we get 
 \[\sqrt{(x-a)^2+(y-b)^2} = \frac{|mx-y+c|}{\sqrt{m^2+1}}.\]Squaring both sides of the equation, we get 
 \[(x-a)^2+(y-b)^2 = \frac{(mx-y+c)^2}{m^2+1}.\]Expanding both sides of the equation gives 
 \[x^2-2ax+a^2+y^2-2by+b^2 = \frac{m^2x^2+y^2+c^2+2mxc-2yc-2mxy}{m^2+1}.\]Multiplying both sides of the equation by $m^2+1$ and rearranging gives 
 \[x^2+2mx+m^2y^2-2((m^2+1)a + mc)x-2((m^2+1)b - c)y +(m^2+1)(a^2+b^2)-c^2.\]Now we can compare to our rotated parabola, $r^2-2rs+s^2+4r = 0$. From this, $-2 = 2m$ or $m = -1$. From here we have a system of three equations:
 \begin{align*} -2((m^2+1)a + mc) &= 4, \\ -2((m^2+1)b - c) &= 0, \\ (m^2+1)(a^2+b^2)-c^2 &= 0. \end{align*}
Plugging in $m=-1$ we get
 \begin{align*} -2(2a-c) &= 4, \\ -2(2b-c) &= 0, \\ 2(a^2+b^2) - c^2 &= 0. \end{align*}
Solving for the first equation, $c = 2+2a$.

 Subtracting the first two equations, $-4a+4b = 4 \implies b = a+1$.

 Plugging into the third equation, $2a^2+2a^2+4a+2 = c^2 \implies 4a^2+4a+2 = c^2$.

 Substituting $c$ in, we get $4a^2+4a+2 = 4a^2+8a+4 \implies 4a+2 = 0 \implies a = -\frac{1}{2}$.

 Now $b = a+1 = \frac{1}{2}$ and $c = 2+2a = 1$.

 This means that the focus of the parabola is $\left(-\frac{1}{2}, \frac{1}{2}\right)$ and the directrix is $y = -x+1$. The maximum value of $r+s$ would lie at the vertex of the parabola, which is the midpoint of the focus and the foot of the focus at the directrix. The line that the vertex and focus lie on is perpendicular to the directrix, so it has slope $1$. It can be written as $y = x+d$ and must go through $\left(-\frac{1}{2}, \frac{1}{2}\right)$ so $d = 1$. This perpendicular line intersects the directrix, so to find the point at which this foot occurs, we set the equation of the lines equal to each other:
 \begin{align*} y &= x+1, \\ y &= -x+1. \end{align*}
Adding, we get $2y = 2$ or $y = 1$ and $x = 0$. The vertex of the parabola is now at the midpoint of $\left(-\frac{1}{2}, \frac{1}{2}\right)$ and $(0, 1)$ which is $\left(-\frac{1}{4}, \frac{3}{4}\right)$.

 Therefore, we have $\boldsymbol{r=-\frac{1}{4}}$ and $\boldsymbol{s=\frac{3}{4}}$.

 



\textbf{Solutions 5.1 and 5.2 Rejoined}

Now that we know the roots of $\tilde{p}(1)$, we can plug in our equation: 
 \[(x-r)(x-s) = \left(1-\left(-\frac{1}{4}\right)\right)\left(1-\frac{3}{4}\right) = \frac{5}{4} \cdot \frac{1}{4} = \frac{5}{16} = \boxed{\textbf{(A) } \frac{5}{16}}.\]~KingRavi

 



\textbf{Solution 6 (Factored Form)}

Let $p(x)=(x-p)(x-q).$ Then, 
 \[p(p(x))=((x-p)(x-q)-p)((x-p)(x-q)-q)=0,\] which means that either $(x-p)(x-q)-p=0$ or $(x-p)(x-q)-q=0$. Both of these equations are quadratics, so $p(p(x))$ has four roots, unless there's a double root. 

 Without loss of generality, let $(x-p)(x-q)-p$ be the expression that produces the double root, so its discriminant is zero. When expanded, 
 \[(x-p)(x-q)-p=x^2-(p+q)x+pq-p=0.\] The value of $x$ is irrelevant to the discriminant, which is $(-(p+q))^2-4(pq-p).$ Setting this equal to zero and simplifing, this equation becomes $p^2-2pq+q^2+4p=0,$ which is a quadratic in $p$.

 
Now, we seek the value of $p$. The previous quadratic is equivalent to $p^2-(2q-4)p+q^2=0$. Using the quadratic formula by having $a=1, b=-(2q-4)$, and $c=q^2,$ we have
 \[p = \frac{2q-4\pm\sqrt{(-(2q-4))^2-4(1(q^2))}}{2(1)} = \frac{2q-4\pm4\sqrt{1-q}}{2}=q-2+2\sqrt{1-q}.\]Our focus is on maximizing $p+q$, so we need the maximum values of $p$ and $q$ respectively (by taking the positive square root). Adding $q,$ we see that
 \[p+q=2\left(q-1+\sqrt{1-q}\right).\]Let $k=\sqrt{1-q}.$ Then, $-k^2=q-1,$ so $p+q=2(-k^2+k).$ 

 
The graph of $-k^2+k$ is a parabola that opens up downwards, and has its maximum value at $y$-value of the vertex. The $x$ coordinate of the vertex is the average of the two roots, which in this case, are $0$ and $1$, so the average of these two is $\frac{1}{2}$. This means that the maximum value of $-k^2+k$ occurs at $k=\frac{1}{2}.$ 

 
Substituting $k=\sqrt{1-q},$ we see that $q=\frac{3}{4}.$ Since $p=q-2+2\sqrt{1-q},$ we can plug $q=\frac{3}{4}$ in, to see that $p=-\frac{1}{4}.$

 
Because the roots of this quadratic are $\frac{3}{4}$ and $-\frac{1}{4},$ our quadratic is $\left(x-\frac{3}{4}\right)\left(x+\frac{1}{4}\right).$ $\tilde{p}(1)=\frac{5}{16}$, so the answer is $\boxed{\textbf{(A) } \frac{5}{16}}$.

 



\textbf{Solution 7 (Vieta’s Formulas)}

Since $\tilde{p}(x)$ is a disrespectful polynomial, it must satisfy $\tilde{p}(\tilde{p}(x))$. Let $\tilde{p}(x)$ have roots $x_1, x_2$. So, the equation $\tilde{p}(\tilde{p}(x))$ becomes two separate equations:
 \begin{align*} x^2+bx+c=x_1, \\ x^2+bx+c=x_2. \end{align*}
Let $r_1,r_2,r_3,r_4$ be the roots of $\tilde{p}(\tilde{p}(x))$. Since there are $3$ real roots, then either one of them is complex, or two of them are the same. It is known that if complex number $w$ is a root to a certain polynomial, then so is $\overline{w}$. So, if one root to $\tilde{p}(\tilde{p}(x))$ is complex, then another root should also be complex — we can't have that. So, two roots are equal; WLOG, assume it is $r_3=r_4$. That means the roots to $\tilde{p}(\tilde{p}(x))$ are $r_1,r_2,r_3,r_3$.

 Now, if the roots of $x^2+bx+c=x_1$ are $r_1$ and $r_3$ and the roots of $x^2+bx+c=x_2$ are $r_2$ and $r_3$, then we have $-b = r_1+r_3 = r_2+r_3$ giving $r_1=r_2$ by \href{https://artofproblemsolving.com/wiki/index.php?title=Vieta\%27s\_Formula}{Vieta's Formula}. That can't happen either. So, $x^2+bx+c=x_1$ must have roots $r_1, r_2$ and $x^2+bx+c=x_2$ must have a double root $r_3$. Again by Vieta's, we have
 \begin{align*} r_1+r_2&=-b, \\ 2r_3&=-b. \end{align*}
That implies $r_2-r_3=r_3-r_1$. So, let $r_1=r_3-k$ and $r_2=r_3+k$. Also, from $2r_3=-b$, we have $r_3=-\frac{b}{2}$. By Vieta's again, we have
 \begin{align*} r_1\cdot r_2 &= c-x_1 = r_1\cdot r_2 = (r_3-k)(r_3+k) = r_3^2-k^2, \\ r_3^2&=c-x_2. \end{align*}
Subtracting yields $x_1-x_2=k^2$. Adding that with $x_1+x_2=-b$ brings us to
 \begin{align*} x_1 &= \frac{k^2-b}{2}, \\ x_2 &= \frac{-k^2-b}{2}. \end{align*}
Thus, $c=x_1\cdot x_2=\frac{b^2-k^4}{4}$. Now, $c-x_2 = r_3^2$ by Vieta's. Plugging in $c, x_2, r_3=-\frac{b}{2}$ gives
 \[\frac{b^2-k^4}{4} + \frac{k^2+b}{2} = \frac{b^2-k^4+2k^2-2b}{4} = \frac{b^2}{4}.\]Simplifying yields $2b=k^4-2k^2=(k^2-1)^2-1 \ge -1$ by \href{https://artofproblemsolving.com/wiki/index.php?title=Trivial\_Inequality}{Trivial Inequality}. Thus, $\text{min}(b) = -\frac{1}{2}$. Since $\tilde{p}(x)$ have the maximum value to the sum of the roots, which is $-b$, we want to minimize $b$. That means that $b$ for $\tilde{p}(x)$ is $\frac{1}{2}$. That means $k^2=1$ and $c=-\frac{3}{16}$, and thus giving us $\tilde{p}(1) = 1+b+c = 1-\frac{1}{2}-\frac{3}{16} = \boxed{\textbf{(A) } \frac{5}{16}}$.

 



\textbf{Solution 8 (Calculus)}

Let $p(x)=(x-r_1)(x-r_2)$. Then 
 \[p(p(x)) = ((x-r_1)(x-r_2)-r_1)((x-r_1)(x-r_2)-r_2)\]As noted in other solutions, in order for $p(p(x))$ to have three (distinct) real roots, one of the roots must be repeated twice.

 For a general polynomial $q(x)$ with an $n$-times repeated root $s$, write $q(x)=(x-s)^n u(x)$, where $u(x)$ is another polynomial. Taking the derivative,
 \begin{align*} q'(x)&= n(x-s)^{n-1} u(x) + (x-s)^n u'(x) \\ &= (x-s)^{n-1} \left[nu(x) + (x-s)u'(x)\right] \end{align*}
This shows that a repeated root / linear factor will appear in the derivative as well.

 Returning to the problem, consider the derivative of $p(p(x))$, which is
 \[[p(p(x))]' = p'(p(x))p'(x) = \left[2(x-r_1)(x-r_2)-r_1-r_2\right] \left[2x-r_1-r_2\right]\]By the observation above, the repeated root of $p(p(x))$ must be a root of this derivative. Therefore, a search through its roots will eventually give information to resolve $r_1$ and $r_2$. Setting the derivative equal to zero yields the candidate equations
 \[2(x-r_1)(x-r_2)-r_1-r_2 = 0 \hspace{.5cm}\text{or}\hspace{.5cm} 2x-r_1-r_2 = 0\]Case 1: Starting with the first equation, if $x_0$ is a solution, then rearrange to obtain
 \[(x_0-r_1)(x_0-r_2) = \frac{r_1+r_2}{2}\]Then
 \begin{align*} p(p(x_0)) &= ((x_0-r_1)(x_0-r_2)-r_1)((x_0-r_1)(x_0-r_2)-r_2) \\ &= \left(\frac{r_1+r_2}{2} - r_1\right)\left(\frac{r_1+r_2}{2} - r_2\right) \\ &= \frac{1}{4}(r_1-r_2)^2 = 0 \end{align*}
This shows that $r_1=r_2$. However, this gives $p(x) = (x-r_1)^2$ and $p(p(x)) = \left[(x-r_1)^2-r_1\right]^2$, which only has at most two distinct real roots. This contradicts the three-root condition given by the problem.

 Case 2: The failure of Case 1 shows that the second equation must contain the repeated root. The solution for the second equation is readily $x_0 = (r_1+r_2)/2$. Then,
 \[p(p(x_0)) = \left[\left(\frac{r_1+r_1}{2}-r_1\right)\left(\frac{r_1+r_2}{2}-r_2\right)-r_1\right]\left[\left(\frac{r_1+r_2}{2}-r_1\right)\left(\frac{r_1+r_2}{2}-r_2\right)-r_2\right] = 0\]Solving for $r_1$ yields four candidate equations:
 \[r_1 = r_2\pm 2\sqrt{-r_2} \hspace{.5cm}\text{or}\hspace{.5cm} r_1 = r_2\pm 2\sqrt{1-r_2} - 2\]Since it is desired to maximise the sum of roots, add $r_2$ to both sides to consider
 \[r_1+r_2 = 2r_2\pm 2\sqrt{-r_2} \hspace{.5cm}\text{or}\hspace{.5cm} r_1+r_2 = 2r_2\pm 2\sqrt{1-r_2} - 2\]For the first family of solutions, the "plus" version is immediately larger, so maximise
 \[f(r_2) = 2r_2 + 2\sqrt{-r_2}\]Taking the derivative and setting equal to zero,
 \[f'(r_2) = 2 - \frac{1}{\sqrt{-r_2}} = 0\]yielding $r_2=-1/4$ and $f(-1/4)=1/2$, which can be verified to be a maximum. Since $f(r_2)=r_1+r_2$, $r_1=3/4$. A similar calculation on the second family of solutions yields the complementary version $r_1=-1/4$ and $r_2=3/4$.

 In summary, this calculation shows that $\tilde{p}(x)=\left(x+\frac{1}{4}\right)\left(x-\frac{3}{4}\right)$. Finally, $\tilde{p}(1)=\left(1+\frac{1}{4}\right)\left(1-\frac{3}{4}\right)=\boxed{\textbf{(A) } \frac{5}{16}}$.

 



\textbf{Solution 9 (Lagrange Multipliers)}

First note that $p(x)$ must have two distinct roots, else $p(p(x))$ will have fewer than $3$ roots. Let these be $a,b.$ Since $p(x)$ has $3$ real roots, and complex roots come in pairs, there must be a double root. So there must be $a_1, a_2$ s.t $p(a_1) = p(a_2) = a$ and $b_1, b_2$ s.t $p(b_1) = p(b_2) = b$ but two of $a_1, a_2, b_1, b_2$ must be the same. WLOG either $a_1 = a_2$ or $a_1 = b_1$. In the latter, case, then, 
 \begin{align*} p(x) &= a + (x - a_1)(x - a_2) \\ p(x) &= b + (x - b_1)(x-b_2) \\ &= b + (x - a_1)(x - b_2) \\ \implies -(a_1 + a_2)x &= -(a_1 + b_2)x \\ \implies a_2 &= b_2 \end{align*}
which implies $a = b$ which is not possible. Hence assume WLOG $a_1 = a_2.$ Then note that
 \begin{align*} p(x) &= (x-a)(x - b) \\  &= a + (x- a_1)^2 \\ -(a + b)x &= -2a_1x, \end{align*}
so
 \begin{align*} p(x) &= a + \left(x - \frac{a + b}{2}\right)^2 \\ x^2 - (a + b)x + ab &= x^2 - (a + b)x + \left(a + \frac{1}{4}(a^2 + 2ab + b^2)\right) \\ 0 &= a^2 - 2ab + b^2 + 4a. \end{align*}
The problem wants us to maximize the sum of the roots, $a + b.$ Thus define $f(a,b) = a + b, g(a,b) = a^2 - 2ab + b^2 + 4a.$ By the method of Lagrange Multipliers, we need to solve,
 \begin{align*} \vec{\nabla} f &= \lambda \vec{\nabla} g \\ g &= 0. \end{align*}
In the first equation, note,
 \begin{align*} \vec{\nabla} f &= \lambda \vec{\nabla} g \\ \langle 1, 1 \rangle &= \lambda \langle 2a - 2b + 4, -2a + 2b \rangle \\ 2a - 2b + 4 &= -2a + 2b \\ a - b &= -1. \end{align*}
Substituting into the second equation, 
 \begin{align*} g(a,b) &= (a-b)^2 + 4a \\ 0 &= 1 + 4a \\ a &= \frac{-1}{4}, b = \frac{3}{4} \\ \tilde{p}(1) &= \left(1 + \frac{1}{4}\right)\left(1 - \frac{3}{4}\right) = \boxed{\textbf{(A) } \dfrac{5}{16}}. \end{align*}
~Aaryabhatta1

 



\textbf{Solution 10 (Solution 8 but Simpler)}

Let $p(x) = (x-r)(x-s)$, having roots $r, s$. WLOG, set $r<s$. The graph of $p(x)$ will have a vertex at $(\frac{r+s}{2}, p(\frac{r+s}{2}))$. Since $p(p(x)) = 0 \iff p(x) = r$ or $p(x) = s$ has exactly 3 distinct solutions, we know that $r = p(\frac{r+s}{2})$ and $s > p(\frac{r+s}{2})$ providing 1 and 2 solutions respectively for a total of 3 (if both $r,s > p(\frac{r+s}{2})$ there will be 4 solutions and if either $r,s<p(\frac{r+s}{2})$ there will be at most 2 solutions). Then,
 \[r = p(\frac{r+s}{2}) = \left(\frac{s-r}{2}\right)\left(\frac{r-s}{2}\right) = -\frac{(s-r)^2}{4}\]
 \[(s-r)^2 = -4r\]
 \[s+r = 2r\pm\sqrt{-4r}\]The negative case clearly doesn't work (giving max $s+r=0$ where $r,s=0$), leaving us with $s+r = 2r+\sqrt{-4r}$. Finding the derivative we see that $s+r$ is maximized when $r = -\frac{1}{4}$, which gives $s = \frac{3}{4}$. Plugging these back into $p(x)$ we see that $\tilde{p}(x)=(x+\frac{1}{4})(x-\frac{3}{4})$. Then $\tilde{p}(1) = \boxed{\textbf{(A) } \dfrac{5}{16}}$.

 



\textbf{Solution 11 (AM-GM Inequality)}

Because $p(x)$ is a quadratic equation with a leading coefficient of 1, let $p(x) = x^2 + ax + b$. Moreover, let $r_1, r_2, r_3$ be the roots of the equation $p(p(x)) = 0$.
 \begin{align*}     p(r_1) &= a \\     p(r_2) &= b \\     p(r_3) &= c     \\\\     p(a) = p(b) &= p(c) = 0 \\     p(x) &= (x - \alpha)(x - \beta) \end{align*}
The following cases are the possible graphs for $p(x)$.
                    \begin{figure}[H]
                        \centering
                        \includegraphics[width=0.5\textwidth]{Solution/AMC_12_12A_Fall_2021_23_Solution_Figure_1.png}
                    \end{figure}
                    The second graph is the case that could provide the maximum sum of the roots.
 \begin{align*}     \therefore \left( x - \frac{\alpha + \beta}{2} \right)^2 + \alpha &= (x - \alpha)(x - \beta) \\     x^2 - (\alpha + \beta)x + \frac{(\alpha + \beta)^2}{4} + \alpha &= x^2 - (\alpha + \beta)x + \alpha\beta \\     \frac{(\alpha + \beta)^2}{4} + \alpha &= \alpha\beta \\     (\alpha + \beta)^2 &= 4\alpha\beta - 4\alpha = 4\alpha(\beta - 1)     \\\\     \sqrt{\alpha(\beta - 1)} &\le \frac{\alpha + (\beta - 1)}{2} \\     &leads\ to \\     \alpha &= \beta - 1.     \\\\     (\alpha + \beta)^2 &= 4\alpha^2 \\     \alpha = \beta\ &or\ 3\alpha + \beta = 0     \\\\     \therefore \alpha = -\frac{1}{4}&, \beta = \frac{3}{4}     \\\\     p(1) = (1 - \alpha)(1 - \beta) &= \frac{5}{4} \cdot \frac{1}{4} = \boxed{\textbf{(A) } \frac{5}{16}} \end{align*}





\textbf{Solution 12: 2 Ways}

In this solution are really 2 solutions: my first one, that I did really fast, but was basically solution 1 except I showed that $\tilde p(\tilde p(x))$ can only have real roots and neither quadratic factor can share any roots. This is too simple, however, so I had to come up with something overkill and unnecessary. So in my second solution, I explicitly list root cases, use derivatives when the optimization could be done with algebra, etc.

 



\textbf{Solution 12.1: The Sane Way}

Let $a$ and $b$ be the roots of $\tilde{p}(x)$. We have that 
 \[\tilde{p}(x) = (x - a)(x - b) \implies \tilde{p}(\tilde{p}(x)) = ((x - a)(x - b) - a)((x - a)(x - b) - b).\]  For this to have $3$ real solutions, one quadratic factor must have a discriminant of $0$ and the other discriminant must be positive. From this, it can quickly be shown that $a$ and $b$ must be real: Assume $a$ and $b$ are nonreal such that $a+b$ and $ab$ are real; then neither quadratic factor can have any real roots because $\tilde p(x)=x^2-(a+b)x+ab$ must be real for all $x\in\mathbb R,$ so for all $x\in\mathbb R$ the equations $\tilde p(x)=a$ and $\tilde p(x)=b$ cannot be satisfied. Then, it quite immediately follows that for the discriminants to differ, $a\neq b$ so $p(x)=a$ and $p(x)=b$ cannot be satisfied simultaneously, which ensures neither quadratic factor shares a root. Thus, without loss of generality:
 \[x^2 - (a + b)x + ab - a=0\text{ for 2 distinct }x\in\mathbb R \implies (a + b)^2 - 4(ab - a) = a^2 + b^2 - 2ab + 4a = 0 \implies (a - b)^2 = -4a \implies a \le 0\]and
 \[x^2 - (a + b)x + ab - b=0\text{ for 1 }x\in\mathbb R \implies (a + b)^2 - 4(ab - b) = a^2 + b^2 - 2ab + 4b > 0 \implies (a - b)^2=-4a > -4b \implies a < b.\] Hence,
 \[b-a=2\sqrt{-a}\implies a+b=2(a+\sqrt{-a}.)\]We want to maximize $a+\sqrt{-a}$ for $a\le 0.$ Substitute $x=\sqrt{-a}$ so we maximize $-x^2+x$ for $x\ge 0.$ We have $-x^2+x=-x(x-1)$ which is maximized at $x=\frac12\implies  a=-\frac14\implies  \max(a+\sqrt{-a})=\frac14.$ Then $\max(a+b)=\frac12,\, a=-\frac14\implies b=\frac34.$ Hence
 \[\tilde{p}(x)=x^2-\frac12 x-\frac3{16}\implies \tilde{p}(1)=\boxed{\textbf{(A) }\frac5{16}}\]~shockfront99

 



\textbf{Solution 12.2: Projectile Motion and General Inefficiency}

Let $q(x)$ be a disrespectful quadratic. It follows that the graph of $q(x)$ is a convex parabola, and thus the graph of $-q(x)$ is a concave parabola. Suppose $q(x)=(x-a)(x-b)$ where $a$ and $b$ are real numbers. Then:
 \[q(q(x))=(q(x)-a)(q(x)-b).\]The graph of $q(q(x))$ can be viewed as essentially graphing $q(t)$ where $t\in \mathrm{ran}(q).$ The following cases yield different amounts of real numbers such that $q(q(x))=0:$

 \textbf{Case 1: 4 real roots} We need the equations
 \[q(x)=a\quad\text{and}\quad q(x)=b\]to have 2 real roots each.

 \textbf{Case 2: 3 real roots} We need the equations 
 \[q(x)=a\quad\text{and}\quad q(x)=b\]to have 2 real roots and 1 real root respectively, or 1 real root and 2 real roots respectively.

 \textbf{Case 3: 2 real roots} We need the equations
 \[q(x)=a\quad\text{and}\quad q(x)=b\]to have 1 real roots each.

 \textbf{Case 4: 1 real root} We need the equations
 \[q(x)=a\quad\text{and}\quad q(x)=b\]to have 1 real root and no real roots respectively, or no real roots and 1 real root respectively.\\

 \textbf{=Case 5: no real roots} We need the equations
 \[q(x)=a\quad\text{and}\quad q(x)=b\]to both have no real roots.

 
Case 2 is what we want. For the $q(x)=a$ and $q(x)=b$ to have different amounts of real solutions, it immediately follows that $a\neq b.$ It can also be quickly shown that to have a disrespectful polynomial, $a$ and $b$ must be real: since any disrespectful quadratic has real coefficients, all real inputs yield real outputs and so the equations $q(x)=a$ and $q(x)=b$ can't have any real solutions unless $a$ and $b$ are real. Let 
 \[\tilde p(x)=(x-\alpha)(x-\beta),\] and fix $\alpha<\beta$ for simplicity. Define the $k$-shift of $g(x)$ as $g(x-k).$ It follows that $-\tilde p(x)$ is concave parabola such that $\max(-\tilde p)>0,$ so $-\tilde p(x)$ is the $\alpha$-shift of the path of a projectile that experiences negligible air resistance launched from a horizontal surface with constant downward acceleration. A projectile launched with speed $u$ at an angle $\theta$ to the horizontal has initial horizontal velocity $u\cos(\theta),$ on which no relevant horizontal drag, friction, or gravitational force effects and thus the projectile has horizontal position 
 \[p_h(t)=ut\cos(\theta)\] after $t$ seconds. Then, the projectile has initial vertical velocity $u\sin(\theta)$ and has a vertical acceleration $g$ downwards and thus a total vertical velocity of $u\sin(\theta)-gt$ and a vertical position of 
 \[p_v(t)=ut\sin(\theta)-\frac12gt^2.\]  We can now see that setting $p^\prime_v(t)=0$ yields that $t=\frac{u\sin(\theta)}{g}$ so letting $H_{\max}=\max(p_v)$ denote the maximum height,
 \[H_{\max}=\frac{u^2\sin^2(\theta)}{g}-\frac{u^2\sin^2(\theta)}{2g}=\frac{u^2\sin^2(\theta)}{2g}.\]Then, to find the range $R$ (horizontal distance the projectile travels) we need $p_v(t)=0$ such that $t\neq0$ so
 \[u\sin(\theta)-\frac{gt}{2}=0\implies t=\frac{2u\sin(\theta)}{g}\implies R=\frac{2u^2\sin(\theta)\cos(\theta)}{g}=\frac{u^2\sin(2\theta)}{g}.\]Now, we can see that modeling $-\tilde p(x)$ as the $\alpha$-shift of such a trajectory, $\beta-\alpha$ is the range of the projectile and so we are maximizing
 \[f(\alpha,\theta,u)=\alpha+(\alpha+R)=2\alpha+\frac{u^2\sin(2\theta)}{g}\]subject to some constraints where $g$ is a positive real constant. Then, the equation $\tilde p(\tilde p(x))=0$ can be rewritten as
 \[-\tilde p(\tilde p(x))=-(\tilde p(x)-\alpha)\left(\tilde p(x)-\alpha-\frac{u^2\sin(2\theta)}{g}\right).\]We now use the constraints of case 2: we either have that $\tilde p(x)-\alpha$ has 2 real roots and $\tilde p(x)-\alpha-\frac{u^2\sin(2\theta)}{g}$ has 1, or vice versa. It can be seen that
 \[-\alpha >-\alpha -\frac{u^2\sin(2\theta)}{g}\](which was immediate from $\alpha<\beta$ anyways) so the vertex of $\tilde p(x)-\alpha-\frac{u^2\sin(2\theta)}{g}$ lies below the vertex of $\tilde p(x)-\alpha.$ Thus $\tilde p(x)-\alpha$ has to have its vertex on the $x$-axis and $\tilde p(x)-\alpha -\frac{u^2\sin(2\theta)}{g}$ has to have its vertex below the $x$-axis. Since $\tilde p(x)-\alpha$ has its vertex on the $x$-axis it must be expressible in the form $(x-h)^2=x^2-2hx+h^2$ for some real $h,$ and be zero at this $h.$ Thus we'd have
 \[0=x^2-\left(2\alpha +\frac{u^2\sin(2\theta)}{g}\right)x+\alpha^2+\frac{\alpha u^2\sin(2\theta)}{g}-\alpha=x^2-2hx+h^2.\]Comparing linear coefficients instantly yields 
 \[h=\alpha+\frac{u^2\sin(2\theta)}{2g}.\]Comparing our constant term and using our expression for $h$ yields
 \[h^2=\alpha^2+\frac{u^2\sin(2\theta)}{g}+\frac{u^4\sin^2(2\theta)}{4g^2}=\alpha^2+\frac{u^2\sin(2\theta)}{g}-\alpha\implies -4\alpha=\frac{u^2\sin^2(2\theta)}{g^2}\implies 2\sqrt{-\alpha} =\frac{u\sin(2\theta)}{g}\quad (\alpha\le0).\]Thus to maximize $f$ we can just maximize
 \[f_1(\alpha)=2\alpha+2\sqrt{-\alpha}\qquad (\alpha\le 0)\]We used common derivatives/integrals for the projectile motion but we're not doing that here for no reason
 \[\frac{\mathrm{d} f}{\mathrm{d} \alpha}=\lim_{h\to 0}\frac{2(\alpha+h)+2\sqrt{-\alpha -h}-2\alpha -2\sqrt{-a} }{h}=2+2\lim_{h\to 0}\frac{\sqrt{-\alpha-h}-\sqrt{-\alpha}}{h}=2+2\lim_{h\to 0}\frac{-h}{h(\sqrt{-\alpha-h}+\sqrt{-\alpha})}=2-\frac{1}{\sqrt {-\alpha}}.\]So
 \[\frac{\mathrm{d} f}{\mathrm{d} \alpha}=0\implies \frac{1}{\sqrt{-\alpha}}=2\implies \sqrt{-\alpha}=\frac12\implies \alpha=-\frac14.\]We then have that $\max(f)=\alpha+\beta=-\frac12+1=\frac12$ so $\beta =\frac34.$ Thus
 \[\tilde p(x)=\left(x+\frac14\right)\left(x-\frac34\right)\implies\tilde p(1)=\frac54\cdot \frac14=\boxed{\textbf{(A) }\frac5{16}}\]~shockfront99

 


	\item (\textit{AMC 12 12A Fall 2021}) Convex quadrilateral $ABCD$ has $AB = 18, \angle{A} = 60^\circ,$ and $\overline{AB} \parallel \overline{CD}.$ In some order, the lengths of the four sides form an arithmetic progression, and side $\overline{AB}$ is a side of maximum length. The length of another side is $a.$ What is the sum of all possible values of $a$?


 $\textbf{(A) } 24 \qquad \textbf{(B) } 42 \qquad \textbf{(C) } 60 \qquad \textbf{(D) } 66 \qquad \textbf{(E) } 84$

 \textbf{Answer: }E



\textbf{Solution 1}

Let $E$ be a point on $\overline{AB}$ such that $BCDE$ is a parallelogram. Suppose that $BC=ED=b, CD=BE=c,$ and $DA=d,$ so $AE=18-c,$ as shown below.
                    \begin{figure}[H]
                        \centering
                        \includegraphics[width=0.5\textwidth]{Solution/AMC_12_12A_Fall_2021_24_Solution_Figure_1.png}
                    \end{figure}
                    We apply the Law of Cosines to $\triangle ADE:$
 \begin{align*} AD^2 + AE^2 - 2\cdot AD\cdot AE\cdot\cos 60^\circ &= DE^2 \\ d^2 + (18-c)^2 - d(18-c) &= b^2 \\ (18-c)^2 - d(18-c) &= b^2 - d^2 \\ (18-c)(18-c-d) &= (b+d)(b-d). \hspace{15mm}(\bigstar) \end{align*}
Let $k$ be the common difference of the arithmetic progression of the side-lengths. It follows that $b,c,$ and $d$ are $18-k, 18-2k,$ and $18-3k,$ in some order. It is clear that $0\leq k<6.$

 If $k=0,$ then $ABCD$ is a rhombus with side-length $18,$ which is valid.

 If $k\neq0,$ then we have six cases:

 
            \begin{enumerate} 
        \item $(b,c,d)=(18-k,18-2k,18-3k)$Note that $(\bigstar)$ becomes $2k(5k-18)=(36-4k)(2k),$ from which $k=6.$ So, this case generates no valid solutions $(b,c,d).$

        \item $(b,c,d)=(18-k,18-3k,18-2k)$Note that $(\bigstar)$ becomes $3k(5k-18)=(36-3k)k,$ from which $k=5.$ So, this case generates $(b,c,d)=(13,3,8).$

        \item $(b,c,d)=(18-2k,18-k,18-3k)$Note that $(\bigstar)$ becomes $k(4k-18)=(36-5k)k,$ from which $k=6.$ So, this case generates no valid solutions $(b,c,d).$

        \item $(b,c,d)=(18-2k,18-3k,18-k)$Note that $(\bigstar)$ becomes $3k(4k-18)=(36-3k)(-k),$ from which $k=2.$ So, this case generates $(b,c,d)=(14,12,16).$

        \item $(b,c,d)=(18-3k,18-k,18-2k)$Note that $(\bigstar)$ becomes $k(3k-18)=(36-5k)(-k),$ from which $k=9.$ So, this case generates no valid solutions $(b,c,d).$

        \item $(b,c,d)=(18-3k,18-2k,18-k)$Note that $(\bigstar)$ becomes $2k(3k-18)=(36-4k)(-2k),$ from which $k=18.$ So, this case generates no valid solutions $(b,c,d).$

            \end{enumerate}
            Together, the sum of all possible values of $a$ is $18+(13+3+8)+(14+12+16)=\boxed{\textbf{(E) } 84}.$




\textbf{Solution 2}

Let $b, c$, and $d$ denote the sides $BC, CD$, and $AD$ respectively.  
                    \begin{figure}[H]
                        \centering
                        \includegraphics[width=0.5\textwidth]{Solution/AMC_12_12A_Fall_2021_24_Solution_Figure_2.png}
                    \end{figure}
                    Since $AB\parallel CD$, we get 
 \[\frac{\sqrt 3}{2}\ d = b\sin\theta \quad \textrm{and}\quad \frac 12 d + c + b\cos\theta = 18.\]Using $b^2\sin^2\theta + b^2\cos^2\theta = b^2$, we eliminate $\theta$ from above to get$(36-2c-d)^2+3d^2=4b^2$, which rearranges to $(36-2c-d)^2-d^2=4(b^2-d^2)$, and, upon factoring, yields
 \begin{align}     (18-c)(18-c-d)=(b+d)(b-d). \end{align}We divide into two cases, depending on whether $c$ is the smallest side.

 If $c$ is not the smallest side then $18-c=\pm (b-d)$. If $c=18$, we get a rhombus of side $18$, so  one possible value is $a=18$. Otherwise, we can cancel the common factor from $(1)$. After rearranging we get
 \[18-c=-b \quad \textrm{or}\quad  18-c=b+2d.\]The first condition is false because $-b< 0 <18-c$; the second condition is false because $b+2d > |b-d| = 18-c$.

 If $c$ is the smallest side, then $18-c = \pm 3(b-d)$. Assuming $c<18$ we can cancel common factors in $(1)$ to get
 \[8b=13d \quad \textrm{or}\quad 8b=7d.\]  The first condition yields the solution $(c,d,b)=(3,8,13)$ and the second condition yields the solution $(c,b,d)=(12,14,16)$.

 Together, the sum of all possible values of $a$ is $18+(3+8+13)+(12+14+16)=\boxed{\textbf{(E) } 84}.$




\textbf{Solution 3}

Denote $x = AD$, $\theta = \angle B$.Hence, $BC = \frac{\sqrt{3}}{2} \cdot \frac{x}{\sin \theta}$, $DC = 18 - \frac{x}{2} - \frac{\sqrt{3}}{2} x \cot \theta$.

 $\textbf{Case 1}$: $DC = AD = BC = AB$.

 This is a rhombus. So each side has length $18$.

 For the following cases, we consider four sides that have distinct lengths.To make their lengths an arithmetic sequence, we must have $\theta \neq 120^\circ$.

 Therefore, in the subsequent analysis, we exclude the solution $\theta = 120^\circ$.

 $\textbf{Case 2}$: $DC < AD < BC < AB$.

 Because the lengths of these sides form an arithmetic sequence, we have the following system of equations:
 \[ AB - BC = BC - AD = AD - DC . \]

 Hence,
 \begin{eqnarray*} & 18 - \frac{\sqrt{3}}{2}\cdot\frac{x}{\sin \theta} = \frac{\sqrt{3}}{2}\cdot\frac{x}{\sin \theta} - x = x - \left( 18 - \frac{x}{2} - \frac{\sqrt{3}}{2} x \cot \theta \right) . & \end{eqnarray*}

 By solving this system of equations, we get $\left( \cos \theta , \sin \theta , x\right) = \left( \frac{11}{13} , \frac{4 \sqrt{3}}{13} , 8 \right)$.

 Thus, in this case, $DC = 3$, $AD = 8$, $BC = 13$.

 $\textbf{Case 3}$: $DC < BC < AD < AB$.

 Because the lengths of these sides form an arithmetic sequence, we have the following system of equations:
 \[ AB - AD = AD - BC = BC - DC . \]

 Hence,
 \begin{eqnarray*} & 18 - x = x - \frac{\sqrt{3}}{2}\cdot\frac{x}{\sin \theta} = \frac{\sqrt{3}}{2}\cdot\frac{x}{\sin \theta}  - \left( 18 - \frac{x}{2} - \frac{\sqrt{3}}{2} x \cot \theta \right) . & \end{eqnarray*}

 By solving this system of equations, we get $\left( \cos \theta , \sin \theta , x\right) = \left( - \frac{1}{7} , \frac{4 \sqrt{3}}{7} , 16 \right)$.

 Thus, in this case, $DC = 12$, $AD = 16$, $BC = 14$.

 $\textbf{Case 4}$: $BC < CD < AD < AB$.

 By doing the similar analysis, we can show there is no solution in this case.

 $\textbf{Case 5}$: $BC < AD < CD < AB$.

 By doing the similar analysis, we can show there is no solution in this case.

 $\textbf{Case 6}$: $AD < CD < BC < AB$.

 By doing the similar analysis, we can show there is no solution in this case.

 $\textbf{Case 7}$: $AD < BC < CD < AB$.

 By doing the similar analysis, we can show there is no solution in this case.

 Therefore, the sum of all possible values of $a$ is
 \begin{align*} 18 + \left( 3 + 8 + 13 \right) + \left( 12 + 14 + 16 \right) & = 84 . \end{align*}


 Therefore, the answer is $\boxed{\textbf{(E) } 84}$.

 


 


	\item (\textit{AMC 12 12A Fall 2021}) Let $m\ge 5$ be an odd integer, and let $D(m)$ denote the number of quadruples $(a_1, a_2, a_3, a_4)$ of distinct integers with $1\le a_i \le m$ for all $i$ such that $m$ divides $a_1+a_2+a_3+a_4$. There is a polynomial

 \[q(x) = c_3x^3+c_2x^2+c_1x+c_0\]such that $D(m) = q(m)$ for all odd integers $m\ge 5$. What is $c_1?$
$\textbf{(A)}\ {-}6\qquad\textbf{(B)}\ {-}1\qquad\textbf{(C)}\ 4\qquad\textbf{(D)}\ 6\qquad\textbf{(E)}\ 11$

 \textbf{Answer: }E



\textbf{Solution 1 (Complete Residue System)}

For a fixed value of $m,$ there is a total of $m(m-1)(m-2)(m-3)$ possible ordered quadruples $(a_1, a_2, a_3, a_4).$

 Let $S=a_1+a_2+a_3+a_4.$ We claim that exactly $\frac1m$ of these $m(m-1)(m-2)(m-3)$ ordered quadruples satisfy that $m$ divides $S:$

 Since $\gcd(m,4)=1,$ we conclude that 
 \[\{k+4(0),k+4(1),k+4(2),\ldots,k+4(m-1)\}\] is the complete residue system modulo $m$ for all integers $k.$

 Given any ordered quadruple $(a'_1, a'_2, a'_3, a'_4)$ in modulo $m,$ it follows that exactly one of these $m$ ordered quadruples has sum $0$ modulo $m:$
 \[\begin{array}{c|c} & \\ [-2.5ex] \textbf{Ordered Quadruple} & \textbf{Sum Modulo }\boldsymbol{m} \\ [0.5ex] \hline & \\ [-2ex] (a'_1, a'_2, a'_3, a'_4) & S'+4(0) \\ [0.5ex] (a'_1+1, a'_2+1, a'_3+1, a'_4+1) & S'+4(1) \\ [0.5ex] (a'_1+2, a'_2+2, a'_3+2, a'_4+2) & S'+4(2) \\ [0.5ex] \cdots & \cdots \\ [0.5ex] (a'_1+m-1, a'_2+m-1, a'_3+m-1, a'_4+m-1) & S'+4(m-1) \\ [0.5ex] \end{array}\]We conclude that $D(m)=\frac1m\cdot[m(m-1)(m-2)(m-3)]=(m-1)(m-2)(m-3),$ so 
 \[q(x)=(x-1)(x-2)(x-3)=c_3x^3+c_2x^2+c_1x+c_0.\]By Vieta's Formulas, we get $c_1=1\cdot2+1\cdot3+2\cdot3=\boxed{\textbf{(E)}\ 11}.$




\textbf{Solution 2 (Symmetric Congruent Numbers and Interpolation)}

Define
 \[ b_i = \left\{ \begin{array}{ll} a_i & \mbox{ if } 1 \leq a_i \leq \frac{m-1}{2} \\ a_i - m & \mbox{ if } \frac{m-1}{2} + 1 \leq a_i \leq m - 1 \\ 0 & \mbox{ if } a_i = m \end{array} \right.. \]

 Hence, $b_i$ is a one-to-one and onto function of $a_i$, and the range of $b_i$ is $\left\{- \frac{m-1}{2}  , \cdots , \frac{m-1}{2} \right\}$.

 Therefore, to solve this problem, it is equivalent for us to count the number of tuples $\left( b_1 , b_2 , b_3 , b_4 \right)$ that are all distinct and satisfy $m | b_1 + b_2 + b_3 + b_4$.

 Denote by $d \left( m \right)$ the number of such tuples that are also subject to the constraint $b_1 < b_2 < b_3 < b_4$.

 Hence, $D \left( m \right) = 4! d \left( m \right) = 24 d \left( m \right)$.

 We do the following casework analysis to compute $d \left( m \right)$.

 $\textbf{Case 1}$: There is one $0$ in $\left( b_1 , b_2 , b_3 , b_4 \right)$.

 Denote by $d_{1i} \left( m \right)$ the number of tuples with $b_i = 0$.

 By symmetry, $d_{11} \left( m \right)= d_{14} \left( m \right)$ and $d_{12} \left( m \right)= d_{13} \left( m \right)$.

 $\textbf{Case 2}$: There is no $0$ in $\left( b_1 , b_2 , b_3 , b_4 \right)$.

 Denote by $d_{2i} \left( m \right)$ the number of tuples with $i$ positive entries.

 By symmetry, $d_{20} \left( m \right) = d_{24} \left( m \right)$ and $d_{21} \left( m \right) = d_{23} \left( m \right)$.

 Therefore,
 \begin{align*} D \left( m \right) & = 24 d \left( m \right) \\ & = 24 \left( \sum_{i=1}^4 d_{1i} \left( m \right) + \sum_{i=0}^4 d_{2i} \left( m \right) \right) \\ & = 24 \left( 2 d_{11} \left( m \right) + 2 d_{12} \left( m \right) + 2 d_{24} \left( m \right) + 2 d_{23} \left( m \right) + d_{22} \left( m \right) \right) . \end{align*}


 Now, we compute $D \left( m \right)$ for $m = 5 , 7 , 9 , 11$. 

 
$\underline{\textbf{SCENARIO}}$ $m = 5$.

 We have $b_i \in \left\{ - 2 , \cdots , 2 \right\}$.

 $\textbf{Case 1}$
$\textbf{Case 1.1}$: $b_1 = 0$

 We cannot have $3$ distinct positive integers. So $d_{11} \left( 5 \right) = 0$.

 $\textbf{Case 1.2}$: $b_2 = 0$

 Because there are $2$ positive integers, we must have $b_3 = 1$, $b_4 = 2$.Hence, $b_1 = - 3$. However, this is out of the range of $b_i$.Thus, $d_{12} \left( 5 \right) = 0$.

 $\textbf{Case 2}$
$\textbf{Case 2.1}$: $b_1 > 0$

 We cannot have $4$ distinct positive integers. So $d_{24} \left( 5 \right) = 0$.

 $\textbf{Case 2.2}$: $b_1 < 0 < b_2$

 We cannot have $3$ distinct positive integers. So $d_{23} \left( 5 \right) = 0$.

 $\textbf{Case 2.3}$: $b_2 < 0 < b_3$

 The only solution is $\left( b_1 , b_2 , b_3 , b_4 \right) = \left( - 2 , - 1 , 1 , 2 \right)$. So $d_{22} \left( 5 \right) = 1$.

 Therefore, $D \left( 5 \right) = 24$.

 
$\underline{\textbf{SCENARIO}}$ $m = 7$.

 We have $b_i \in \left\{ - 3 , \cdots , 3 \right\}$.

 $\textbf{Case 1}$
$\textbf{Case 1.1}$: $b_1 = 0$

 We have no feasible solution. Thus, $d_{11} \left( 7 \right) = 0$.

 $\textbf{Case 1.2}$: $b_2 = 0$

 The only solution is $\left( b_1 , b_3 , b_4 \right) = \left( - 3 , 1 , 2 \right)$.Thus, $d_{12} \left( 7 \right) = 1$.

 $\textbf{Case 2}$
$\textbf{Case 2.1}$: $b_1 > 0$

 We cannot have $4$ distinct positive integers. So $d_{24} \left( 7 \right) = 0$.

 $\textbf{Case 2.2}$: $b_1 < 0 < b_2$

 To get $3$ distinct positive integers, we have $\left( b_2 , b_3 , b_4 \right) = \left( 1 , 2 , 3 \right)$. This implies $b_1 = - 6$. However, this is out of the range of $b_1$. So $d_{23} \left( 6 \right) = 0$.

 $\textbf{Case 2.3}$: $b_2 < 0 < b_3$

 We have $d_{22} \left( 7 \right) = \binom{3}{2} = 3$.

 Therefore, $D \left( 7 \right) = 24 \cdot 5$.

 
$\underline{\textbf{SCENARIO}}$ $m = 9$.

 We have $b_i \in \left\{ - 4 , \cdots , 4 \right\}$.

 $\textbf{Case 1}$
$\textbf{Case 1.1}$: $b_1 = 0$

 The only solution is $\left( b_2 , b_3 , b_4 \right) = \left(  2, 3, 4 \right)$.Thus, $d_{11} \left( 9 \right) = 1$.

 $\textbf{Case 1.2}$: $b_2 = 0$

 The feasible solutions are $\left( b_1 , b_3 , b_4 \right) = \left( - 3 , 1 , 2 \right)$, $\left( - 4 , 1 , 3 \right)$.Thus, $d_{12} \left( 9 \right) = 2$.

 $\textbf{Case 2}$
$\textbf{Case 2.1}$: $b_1 > 0$

 There is no feasible solution. So $d_{24} \left( 9 \right) = 0$.

 $\textbf{Case 2.2}$: $b_1 < 0 < b_2$

 To get $3$ distinct positive integers, we have $b_2 + b_3 + b_4 \geq 1 + 2 + 3 = 6$. This implies $b_1 = - 6$. However, this is out of the range of $b_1$. So $d_{23} \left( 9 \right) = 0$.

 $\textbf{Case 2.3}$: $b_2 < 0 < b_3$

 We have $d_{22} \left( 9 \right) = 8$.

 Therefore, $D \left( 9 \right) = 24 \cdot 14$.

 
$\underline{\textbf{SCENARIO}}$ $m = 11$.

 We have $b_i \in \left\{ - 5 , \cdots , 5 \right\}$.

 $\textbf{Case 1}$
$\textbf{Case 1.1}$: $b_1 = 0$

 The only solution is $\left( b_2 , b_3 , b_4 \right) = \left(  2, 4, 5 \right)$.Thus, $d_{11} \left( 11 \right) = 1$.

 $\textbf{Case 1.2}$: $b_2 = 0$

 The feasible solutions are $\left( b_1 , b_3 , b_4 \right) = \left( - 3 , 1 , 2 \right)$, $\left( - 4 , 1 , 3 \right)$, $\left( - 5 , 1 , 4 \right)$, $\left( - 5 , 2 , 3 \right)$.Thus, $d_{12} \left( 11 \right) = 4$.

 $\textbf{Case 2}$
$\textbf{Case 2.1}$: $b_1 > 0$

 The only feasible solution is $\left( b_1 , b_2 , b_3 , b_4 \right) = \left(  1 , 2, 3, 5 \right)$. So $d_{24} \left( 11 \right) = 1$.

 $\textbf{Case 2.2}$: $b_1 < 0 < b_2$

 The only feasible solution is $\left( b_1 , b_2 , b_3 , b_4 \right) = \left(  -1 , 3, 4, 5 \right)$.  So $d_{23} \left( 11 \right) = 1$.

 $\textbf{Case 2.3}$: $b_2 < 0 < b_3$

 We have $d_{22} \left( 11 \right) = 16$.

 Therefore, $D \left( 11 \right) = 24 \cdot 30$.

 We know that $q \left( m \right) = D \left( m \right)$ for odd $m \geq 5$.

 Plugging $m = 5, 7, 9, 11$ into this equation, we get
 \begin{align*} c_3 5^3 + c_2 5^2 + c_1 5 + c_0 & = 24 \cdot 1 && (1.1) \\ c_3 7^3 + c_2 7^2 + c_1 7 + c_0 & = 24 \cdot 5 && (1.2) \\ c_3 9^3 + c_2 9^2 + c_1 9 + c_0 & = 24 \cdot 14 && (1.3) \\ c_3 11^3 + c_2 11^2 + c_1 11 + c_0 & = 24 \cdot 30 && (1.4) \end{align*}


 Now, we solve this system of equations.Taking $\frac{(1.2)-(1.1)}{2}$, $\frac{(1.3)-(1.2)}{2}$, $\frac{(1.4)-(1.2)}{2}$, we get
 \begin{align*} c_3 109 + c_2 12 + c_1 & = 48 && (2.1) \\ c_3 193 + c_2 16 + c_1 & = 108 && (2.2) \\ c_3 301 + c_2 20 + c_1 & = 192 && (2.3) \end{align*}


 Taking $\frac{(2.2)-(2.1)}{4}$, $\frac{(2.3)-(2.2)}{4}$, we get
 \begin{align*} c_3 21 + c_2  & = 15 && (3.1) \\ c_3 27 + c_2 & = 21 && (3.2) \end{align*}


 Taking $\frac{(3.2)-(3.1)}{6}$, we get $c_3 = 1$.

 Plugging $c_3$ into Equation (3.1), we get $c_2 = - 6$.

 Plugging $c_2$ and $c_3$ into Equation (2.1), we get $c_1 = 11$.

 Therefore, the answer is $\boxed{\textbf{(E) }11}$.

 



\textbf{Solution 3}

As before, note that we have $m(m-1)(m-2)$ numbers we can choose as $a,b,c.$ From here, there is exactly one possible value of $1 \leq d \leq m$ that could make $S=a+b+c+d$ divisible by $m.$ However, there is a $\frac{3}{m}$ chance that this value of $d$ has already been chosen as $a,b$ or $c$. Thus our polynomial is $m(m-1)(m-2)\left(1-\frac{3}{m}\right)=m(m-1)(m-2)\left(\frac{m-3}{m}\right)=(m-1)(m-2)(m-3)$. By Vieta's, $c_1 = 2+3+6=\boxed{\textbf{(E) } 11}$.

 



\textbf{Solution 4}

We can solve this by directly counting quadruples and deriving a closed formula for $D(m)$.

 For an odd integer $m \geq 5$, we need to count ordered quadruples $(a_1,a_2,a_3,a_4)$ of distinct integers where $1 \leq a_i \leq m$ and $m$ divides their sum.

 Consider how many total ordered quadruples of distinct integers exist in the range $[1,m]$:- We choose 4 distinct numbers from $m$ possible values: $\binom{m}{4}$- We arrange these 4 numbers in order: $4!$

 Therefore, total quadruples = $\binom{m}{4} \cdot 4! = \frac{m!}{(m-4)!}$

 Now, these quadruples can be partitioned into $m$ groups based on the remainder of their sum when divided by $m$. By symmetry, these remainders are uniformly distributed modulo $m$.

 Therefore, exactly $\frac{1}{m}$ of all quadruples have a sum divisible by $m$:

 
 \[D(m) = \frac{1}{m} \cdot \frac{m!}{(m-4)!} = \frac{m(m-1)(m-2)(m-3)}{m} = (m-1)(m-2)(m-3)\]

 Expanding this expression:
 \[(m-1)(m-2)(m-3) = m^3 - 6m^2 + 11m - 6\]

 Comparing with $q(x) = c_3x^3 + c_2x^2 + c_1x + c_0$, we identify:$c_3 = 1$, $c_2 = -6$, $c_1 = 11$, $c_0 = -6$

 Therefore, $c_1 = 11$, and the answer is $\boxed{(E) 11}$.

 



\textbf{Solution 5 (Principle of Inclusion-Exclusion)}

If we do not have to worry about whether or not the numbers are distinct, then there would be $m^3$ possible quadruples. This is because the first $3$ numbers can be anything, and then the last number is uniquely chosen so that the sum is a multiple of $m$.

 Now, we must subtract the quadruples with at least two of the same number. There are $\binom{4}{2}$ ways to choose the positions for the two identical numbers, $m$ ways to choose the number for these two positions, and $m$ ways to choose a number for another position. The last position is uniquely determined as before. Thus, our expression becomes $m^3-\binom{4}{2}\cdot m\cdot m = m^3 - 6m^2$.

 However, we have subtracted the cases in which there are three identical numbers too many times. Each such case was subtracted $\binom{3}{2}=3$ times, and there are $\binom{4}{3}\cdot m=4m$ such cases. Thus, we need to add back $8m$.

 We also subtracted the cases in which there are two pairs of two identical numbers too many times. There are $\frac{\binom{4}{2}m}{2}=3m$ such cases. This is because we first choose the two positions for the first two identical numbers, then multiply by $m$ to determine which number goes in these positions. The last two positions can be determined uniquely because $m$ is odd. We divide by $2$ because we counted each case twice. In our original counting, we subtracted each such case twice, so we add back $3m$. In total, our expression is now $m^3-6m^2+8m+3m=m^3-6m^2+11m$.

 We can stop here, because we know that the number of times $(5,5,5,5)$ was overcounted is just a constant. Thus, the answer is $\boxed{(E) 11}$.

 


\end{enumerate}